\chapter{世界观基本规律的理论介绍}\label{chap:theory}
本章主要介绍OC所在的世界观的基本规律的介绍。涵盖以下内容：
\begin{introduction}
    \item \hyperref[sec:usefulfunction]{能力函数理论}
    \item \hyperref[sec:worldtheory]{世界理论}
    \item \hyperref[sec:elementtheory]{元素理论}
    \item \hyperref[sec:octheory]{OC 变换理论}
    \item \hyperref[sec:ocimage]{OC的数字图像生成}
\end{introduction}

这些理论基于数学与逻辑学的基本原理，并运用计算机科学的部分概念进行描述
。通过这些理论，可以更好地理解\wkn{}的世界观的运作机制。

读者在阅读本章时，需要具备一定的数学基础，尤其是微积分、逻辑学等基本知识。此外，还需进行一定的计算机科学的了解，例如面向对象编程、基础数据结构、信号处理、数字电路等。

\iftoggle{appdx}{。
    本书的附录部分提供了相关的数学与计算机科学基础知识，供读者参考。具体内容可见：
    \begin{itemize}
        \item \textbf{附录A：数学基础}（见\hyperref[ch:math]{“世界观的数学基础”}），涵盖向量/矢量、高等代数、解析几何、微积分、复变函数、概率统计等内容。
        \item \textbf{附录B：计算机科学基础}（见\hyperref[ch:computer]{“世界观的计算机基础”}），涵盖面向对象设计、基础数据结构、离散数学的代数系统、信号处理、数字电路等内容。
    \end{itemize}
}{
    限于篇幅，本章不包含数学与计算机科学的基础知识介绍，读者可参考相关教材进行补充学习。
}

\section{能力函数理论}\label{sec:usefulfunction}
这一节主要介绍能力函数理论。能力函数是二次元世界中角色能力的数学抽象，是理解和设计角色能力的基础。
\begin{introduction}
    \item \hyperref[subsec:functionandusefulfunction]{能力函数的定义}
    \item \hyperref[subsec:discreteandcontinuous]{能力函数的两种基本类型}
    \item \hyperref[subsec:spell]{符咒}
    \item \hyperref[subsec:gradient]{向量微分符咒}
\end{introduction}

\subsection{能力函数的定义}\label{subsec:functionandusefulfunction}

在二次元世界的数学模型中，角色的特殊能力可以被抽象为一种映射关系，我们称之为\textbf{能力函数（Ability Function）}。此函数严格描述了能力如何将输入（作用目标或条件）转化为输出（作用效果）。

\begin{definition}[能力函数]
    设 $\mathcal{X}$ 为能力的\textbf{输入空间 (Input Space)}，包含了所有可能的能力作用目标、初始条件或触发状态。设 $\mathcal{Y}$ 为能力的\textbf{输出空间 (Output Space)}，包含了所有可能的能力作用效果。

    一个\textbf{能力函数 (Ability Function)} 是一个映射：
    \begin{equation}
        F: \mathcal{X} \to \mathcal{Y}
    \end{equation}
    它将每个输入 $x \in \mathcal{X}$ 映射到唯一的输出 $y = F(x) \in \mathcal{Y}$。

    此时称：
    \begin{itemize}
        \item $x$ 为能力的\textbf{输入}或\textbf{原像}。
        \item $y = F(x)$ 为能力在输入 $x$ 下的\textbf{输出}或\textbf{像}。
        \item 映射 $F$ 的具体形式，体现了角色能力的使用规则与使用效果。
    \end{itemize}
\end{definition}

能力函数具有以下重要特征：
\begin{enumerate}
    \item \textbf{输入与输出空间的广泛性}：输入空间 $\mathcal{X}$ 和输出空间 $\mathcal{Y}$ 不限于传统的数集，可以是：
          \begin{itemize}
              \item 角色、物体等实体对象的集合。
              \item HP、MP、状态等属性值的集合。
              \item 空间坐标、时间区间等连续范畴。
              \item 向量、矩阵，或其他抽象的数学结构。
          \end{itemize}
    \item \textbf{世界观的约束}：同一世界观下的所有能力函数必须遵守该世界设定的基本规律。
    \item \textbf{确定性}：对于确定的输入 $x$，能力函数必须产生确定的输出 $F(x)$。
\end{enumerate}

\subsection{能力函数的两种类型}\label{subsec:discreteandcontinuous}

根据输入 $x$ 的数学性质及其被处理的方式，能力函数在实践中表现为两种基本的作用模式，进而分为两种基本类型：离散型能力函数和连续型能力函数。
\subsubsection{离散型能力函数}

当能力函数的输入是离散的、可数的对象时，其作用模式是离散的。

\begin{definition}[离散型能力函数]
    若能力函数 $F$ 的输入空间 $\mathcal{X}$ 是一个离散集合（有限或可数无限），则称 $F$ 为\textbf{离散型能力函数 (Discrete Ability Function)} 运行。其效果通过直接将函数作用于每个输入元素来实现：
    \begin{equation}
        \mathcal{Y} = F(\mathcal{X}) = \{F(x) \mid x \in \mathcal{X}\}
    \end{equation}
    若本次能力发动仅针对输入集合 $X = \{x_1, x_2, \ldots, x_n\} \subseteq \mathcal{X}$，则对应的输出集合\(\mathcal{Y} = \{y_1, y_2, \ldots, y_n\}\)为：
    \begin{equation}
        \forall x_i \in X, \quad y_i = F(x_i)
    \end{equation}
    其中 $X \subseteq \mathcal{X}$ 是本次能力发动所针对的具体目标集合。
\end{definition}

\begin{example}[单体治愈术]
    角色A的治愈术以单一角色为目标，恢复其50点HP。定义输入空间 $\mathcal{X}$ 为所有角色的HP状态集合，输出空间 $\mathcal{Y}$ 同样为HP状态集合。能力函数为：
    \begin{equation}
        F_{\text{heal}}(\text{HP}) = \text{HP} + 50
    \end{equation}
    若对目标 $x$（其当前HP为 $\text{HP}_x$）使用该能力，则效果为 $\text{HP}'_x = F_{\text{heal}}(\text{HP}_x)$。
\end{example}

\subsubsection{连续型能力函数}

当能力函数的输入是一个连续空间的元素，且其效果需要通过“累积”来计算时，其作用模式是连续的。

\begin{definition}[连续型能力函数]
    若能力函数 $F$ 的输入空间是一个连续域 $\Omega = \mathcal{X}$（例如一片空间区域或一段时间区间），并且总效果由定义在 $\Omega$ 上的\textbf{效果密度函数 (Effect Density Function)} $f(\omega)$ 积分得到，则称 $F$ 为\textbf{连续型能力函数 (Continuous Ability Function)} 。其总效果集合 $\mathcal{Y}$ 表示为：
    \begin{equation}
        \mathcal{Y} = F(\Omega) = \int_{\Omega} f(\omega)  \md\omega
    \end{equation}
    其中 $\omega$ 是域 $\Omega$ 上的微元，$f: \Omega \to \mathbb{R}$ 是密度函数，它本身可以看作一个局部的、点上的能力函数。
\end{definition}

\begin{remark}
    从形式上看，连续性型能力函数与离散型能力函数的表达式并无区别，都是\(\mathcal{Y}=F(\mathcal{X})\)。但其计算方式和应用场景有本质区别：
    \begin{itemize}
        \item 离散型能力函数：表达式为\(\mathcal{Y}=F(\mathcal{X})\)可以直接对每个离散输入进行函数求值，适用于单体或有限目标的能力作用。
        \item 连续型能力函数：虽然表达式为\(\mathcal{Y}=F(\Omega)=F(\mathcal{X})\)，但无法直接进行求值，而是通过积分对连续域内的效果密度进行累积，适用于范围性、区域性或时间段性的能力作用。
    \end{itemize}
\end{remark}

\begin{example}[范围火球术]
    角色B释放火球术，对圆形区域 $\Omega = \{\vec{r} \in \mathbb{R}^2 : \|\vec{r}\| \leq R\}$ 造成伤害。伤害密度在中心最高，向外衰减，定义为 $f(\vec{r}) = 100 e^{-\|\vec{r}\|^2}$。则该能力的总伤害输出为：
    \begin{align*}
        \text{Damage} & = F_{\text{fireball}}(\Omega) = \iint_{\Omega} f(\vec{r})  \md S \\
                      & = \iint_{\Omega} 100 \me^{-(x^2 + y^2)}  \md x \md y
    \end{align*}
    由于被积函数是圆对称的，我们采用极坐标变换求解，令 $x = \rho \cos\theta,  y = \rho \sin\theta$，则 $\md S = \rho  \md \rho \md \theta$，积分区域 $\Omega$ 变为 $\{(\rho, \theta) \mid 0 \le \rho \le R, 0 \le \theta \le 2\pi\}$。
    \begin{align*}
        \text{Damage} & = \int_{0}^{2\pi} \int_{0}^{R} 100 \me^{-\rho^2} \cdot \rho  \md \rho \md \theta \\
                      & = 100 \int_{0}^{2\pi} \md \theta \int_{0}^{R} \rho \me^{-\rho^2} \md \rho        \\
                      & = 100 \cdot 2\pi \int_{0}^{R} \rho \me^{-\rho^2} \md \rho
    \end{align*}
    对于内层积分，令 $u = \rho^2$，则 $\md u = 2\rho \md \rho$，即 $\rho \md \rho = \frac{\md u}{2}$。
    \begin{align*}
        \int \rho \me^{-\rho^2} \md \rho         & = \int \me^{-u} \cdot \frac{\md u}{2} = -\frac{1}{2}\me^{-u} + C = -\frac{1}{2}\me^{-\rho^2} + C                                  \\
        \int_{0}^{R} \rho \me^{-\rho^2} \md \rho & = \left[ -\frac{1}{2}\me^{-\rho^2} \right]_{0}^{R} = -\frac{1}{2}\me^{-R^2} - (-\frac{1}{2}\me^{0}) = \frac{1}{2}(1 - \me^{-R^2})
    \end{align*}
    代入总伤害公式：
    \begin{align*}
        \text{Damage} & = 100 \cdot 2\pi \cdot \frac{1}{2}(1 - \me^{-R^2}) \\
                      & = 100\pi (1 - \me^{-R^2})
    \end{align*}
    因此，该范围火球术的总伤害函数为：
    \begin{equation}
        F_{\text{fireball}}(\Omega) = 100\pi (1 - \me^{-R^2})
    \end{equation}

    这是一个典型的连续型能力函数。值得注意的是，该能力对区域内位于点 $\vec{r}_p$ 的单个目标造成的伤害，是密度函数在该点的取值 $f(\vec{r}_p) = 100 \me^{-\|\vec{r}_p\|^2}$，这体现了连续型能力函数与离散型能力函数在计算上的关联与区别。
\end{example}

\subsection{符咒}\label{subsec:spell}
\textbf{符咒（Spell）}是世界观中一种特殊的能力函数算子，它通过改变能力函数本身的映射规则来创造新的能力。

\begin{definition}[符咒]
    设 $\mathscr{F}$ 为某世界观下所有能力函数构成的集合。一个\textbf{符咒（Spell）} $\mathcal{O}$ 是一个从能力函数集到其自身的映射：
    \begin{equation}
        \mathcal{O}: \mathscr{F} \to \mathscr{F}
    \end{equation}
    它将一个能力函数 $F \in \mathscr{F}$ 映射为另一个新的能力函数 $\mathcal{O}F \in \mathscr{F}$。

    对于任意作用对象 $x \in \mathcal{X}_F$，新能力函数的效果由下式确定：
    \begin{equation}
        (\mathcal{O}F)(x)
    \end{equation}
    符咒 $\mathcal{O}$ 改变了能力的作用规则，而不仅仅是其输出值。
\end{definition}

符咒与简单函数复合的关键区别在于：符咒操作的是函数本身这一映射规则，而非具体的输入输出值。常见的符咒类型包括：

\begin{itemize}
    \item \textbf{线性强化符咒} $\mathcal{E}_k$：重新定义能力函数的输出为原输出的 $k$ 倍
          \begin{equation}
              (\mathcal{E}_k F)(x) = k \cdot F(x)
          \end{equation}

    \item \textbf{指数衰减符咒} $\mathcal{D}_{\lambda}$：使能力函数效果随空间或时间坐标指数衰减。设 $F(\vec{r})$ 为空间依赖的能力函数，则该符咒作用后为：
          \begin{equation}
              (\mathcal{D}_{\lambda} F)(\bm{r}) = F(\bm{r}) \cdot \me^{-\lambda \|\bm{r}\|}
          \end{equation}

    \item \textbf{作用域限制符咒} $\mathcal{R}_S$：将能力函数的有效作用域限制在区域 $S$ 内
          \begin{equation}
              (\mathcal{R}_S F)(x) = \begin{cases}
                  F(x) & \text{若 } x \in S    \\
                  0    & \text{若 } x \notin S
              \end{cases}
          \end{equation}

\end{itemize}

符咒的复合运算 $\mathcal{O}_2 \mathcal{O}_1$ 定义为：
\begin{equation}
    (\mathcal{O}_2 \mathcal{O}_1)F = \mathcal{O}_2(\mathcal{O}_1 F)
\end{equation}
符咒的复合通常不满足交换律，即 $\mathcal{O}_2 \mathcal{O}_1 \neq \mathcal{O}_1 \mathcal{O}_2$。

符咒的合理组合与运用是\wkn 的世界观中角色能力开发的核心机制。

\subsection{向量微分符咒}\label{subsec:gradient}

\textbf{向量微分符咒（Nabla Spell）}是\wkn 的世界观中最为强大的空间变化符咒之一，一般写作$\nabla$。该符咒可以同时操控空间中的多种变化率，根据不同的组合方式实现高斯公式、格林公式等强大效果。

\subsubsection{向量微分符咒的概念}

\begin{definition}[向量微分符咒]
    在三维空间中，向量微分符咒$\nabla$定义为：
    \begin{equation}
        \nabla = \left(\dfrac{\partial}{\partial x}, \dfrac{\partial}{\partial y}, \dfrac{\partial}{\partial z}\right)
    \end{equation}
    该符咒可以作用于标量能力函数$\varphi$或向量能力函数$\bm{F}=(F_x,F_y,F_z)$，产生不同效果：
    \begin{align*}
        \text{梯度符咒} & : \nabla\varphi = \left(\dfrac{\partial\varphi}{\partial x}, \dfrac{\partial\varphi}{\partial y}, \dfrac{\partial\varphi}{\partial z}\right)                                                                                                      \\
        \text{散度符咒} & : \nabla\cdot\bm{F} = \dfrac{\partial F_x}{\partial x} + \dfrac{\partial F_y}{\partial y} + \dfrac{\partial F_z}{\partial z}                                                                                                                      \\
        \text{旋度符咒} & : \nabla\times\bm{F} = \left(\dfrac{\partial F_z}{\partial y} - \dfrac{\partial F_y}{\partial z}, \dfrac{\partial F_x}{\partial z} - \dfrac{\partial F_z}{\partial x}, \dfrac{\partial F_y}{\partial x} - \dfrac{\partial F_x}{\partial y}\right)
    \end{align*}
    而在二维平面上，$\nabla$符咒可以类似地定义为：
    \begin{equation}
        \nabla_{\text{2D}} = \left(\dfrac{\partial}{\partial x}, \dfrac{\partial}{\partial y}\right),\,
        \nabla_{\text{2D}}\times \bm{F}=\dfrac{\partial F_y}{\partial x} - \dfrac{\partial F_x}{\partial y}
    \end{equation}
\end{definition}

\begin{remark}
    当一个连续性能力函数作用的时候，会形成一个场。这个场又称作\textbf{能力场（Ability Field）}。
\end{remark}

\subsubsection{向量微分符咒的应用}
向量微分符咒$\nabla$可以产生强大的效果。主要的效果就是\textbf{作用域的升维}。

\begin{theorem}[高斯符咒公式]
    当$\nabla$符咒与闭合曲面积分$\displaystyle\oiint$组合时，可实现\textbf{作用域从曲面域到空间域的变化}：
    \begin{equation}
        \oiint_{\partial V} \bm{F}\cdot\mathrm{d}\bm{S}=\iiint_V (\nabla\cdot\bm{F}) \,\md V
    \end{equation}
    此时作用范围扩大，由空间的包裹曲面域变化到了空间域。
\end{theorem}

\begin{theorem}[格林符咒公式]
    在二维平面上，$\nabla_{\text{2D}}$符咒与闭合曲线积分$\dsy\oint$组合：
    \begin{equation}
        \oint_C (L\mathrm{d}x + M\mathrm{d}y) = \iint_D \left(\dfrac{\partial M}{\partial x} - \dfrac{\partial L}{\partial y}\right) \mathrm{d}x\mathrm{d}y
    \end{equation}
    此时实现了作用范围由曲线到曲线包围的平面的变化。
\end{theorem}

\begin{theorem}[斯托克斯符咒公式]
    对于空间曲线，$\nabla$符咒与闭合曲线积分$\dsy\oint$组合：
    \begin{equation}
        \oint_C \bm{F}\cdot\mathrm{d}\bm{r} = \iint_S (\nabla\times\bm{F})\cdot\mathrm{d}\bm{S}
    \end{equation}
    此时作用范围由空间曲线转化为空间曲面。
\end{theorem}

\begin{example}[向量微分符咒综合应用]

    「空间温度操控者」绫的能力函数为：
    \begin{itemize}
        \item 温度感知能力函数：$\varphi(x,y,z)=2x^{2}y+yz^{3}+3z$；
        \item 温度流动能力函数：$\bm{F}(x,y,z)=(xy^{2},\,z\sin y,\,x^{2}z)$。
    \end{itemize}

    \begin{enumerate}
        \item 求点$P(1,0,2)$处的梯度$\nabla\varphi$，并说明该向量如何指示温度变化率最强的方向。

        \item
              绫在半径为$2$的球面$S:\,x^{2}+y^{2}+z^{2}=4$上施展温度流动能力组合。\\
              计算绫仅作用于该球面时的结果
              \begin{equation}
                  \oiint_{S}\bm{F}\cdot\md\bm{S}
              \end{equation}
              和绫作用于球体$x^{2}+y^{2}+z^{2}\leq 4$时的结果
              \begin{equation}
                  \iiint_{V}(\nabla\cdot\bm{F})\,\md V
              \end{equation}
              验证两侧数值相等，以此说明：向量微分符咒能借助闭合曲面把「曲面作用域」升格为「体积作用域」。
        \item
              设在平面$z=1$内，绫取逆时针方向的圆$C:\,x^{2}+y^{2}=9,\,z=1$作为闭合曲线，并令曲面$S$为同一圆盘
              \begin{equation}
                  S:\,x^{2}+y^{2}\leq 9,\,z=1\quad(\text{法向量朝上})。
              \end{equation}
              绫只会在曲线上施展能力\(F\)，但她更想作用于整个曲面。\\
              请帮助绫完成整个曲面的能力作用。
    \end{enumerate}

    \solution
    \begin{enumerate}
        \item $\nabla\varphi=(4xy,\,2x^{2}+z^{3},\,3yz^{2}+3)$，代入$P(1,0,2)$得
              \begin{equation}
                  \nabla\varphi\big|_{P}=(0,\,10,\,15)。
              \end{equation}
              该向量指向温度上升最剧烈的方向，其模长给出最大变化率$\sqrt{10^{2}+15^{2}}=\sqrt{325}\,\mathrm{K/m}$。

        \item 仅作用于球面，有
              \begin{equation}
                  \oiint_{S}\bm{F}\cdot\md\bm{S} = \oiint_{S}(xy^{2},\,z\sin y,\,x^{2}z)\cdot\md\bm{S}
              \end{equation}
              法向量$\md\bm{S}=\left(\dfrac{x}{\sqrt{x^{2}+y^{2}+z^{2}}},\,\dfrac{y}{\sqrt{x^{2}+y^{2}+z^{2}}},\,\dfrac{z}{\sqrt{x^{2}+y^{2}+z^{2}}}\right)\md S$，代入球面方程
              \begin{equation}
                  \md\bm{S}=\left(\dfrac{x}{2},\,\dfrac{y}{2},\,\dfrac{z}{2}\right)\md S。
              \end{equation}
              故有
              \begin{equation}
                  \oiint_{S}\bm{F}\cdot\md\bm{S}=\oiint_{S}\left(\dfrac{x^{2}y^{2}}{2}+\,\dfrac{yz\sin y}{2}+\,\dfrac{x^{2}z^{2}}{2}\right)\md S。
              \end{equation}
              由对称性可知，$\dsy\oiint_{S}yz\sin y\md S=0$。\\
              则：
              \begin{equation}
                  \oiint_S \bm{F} \cdot \md\bm{S}=\oiint_{S}\left(\dfrac{x^{2}y^{2}}{2}+\,\dfrac{x^{2}z^{2}}{2}\right)\md S
              \end{equation}
              又有\(y^2+z^2=4-x^2\)，则
              \begin{equation}
                  \oiint_S \bm{F} \cdot \md\bm{S} = \frac12 \oiint_S x^2(4 - x^2) \, dS
              \end{equation}

              拆分积分：
              \begin{equation}
                  = \frac12 \left[ 4 \oiint_S x^2 \, \md S - \oiint_S x^4 \, \md S \right]
              \end{equation}

              计算第一项：由对称性，
              \begin{equation}
                  \oiint_S x^2 \, \md S = \oiint_S y^2 \, \md S = \oiint_S z^2 \, \md S
              \end{equation}
              且
              \begin{equation}
                  \oiint_S (x^2 + y^2 + z^2) \, \md S = \oiint_S 4 \, \md S = 4 \times 16\pi = 64\pi
              \end{equation}
              所以
              \begin{equation}
                  3 \oiint_S x^2 \, \md S = 64\pi \quad\Rightarrow\quad \oiint_S x^2 \, \md S = \frac{64\pi}{3}
              \end{equation}

              计算第二项：用球坐标
              \begin{equation}
                  x = 2\sin\theta\cos\phi,\quad \md S = 4\sin\theta \, \md\theta \, \md\phi
              \end{equation}
              \begin{equation}
                  \oiint_S x^4 \, \md S = \int_0^{2\pi} \int_0^{\pi} (2\sin\theta\cos\phi)^4 \cdot 4\sin\theta \, \md\theta \, \md\phi
              \end{equation}
              \begin{equation}
                  = 64 \int_0^{2\pi} \cos^4\phi \, \md\phi \cdot \int_0^{\pi} \sin^5\theta \, \md\theta
              \end{equation}
              已知
              \begin{equation}
                  \int_0^{2\pi} \cos^4\phi \, \md\phi = \frac{3\pi}{4},\quad
                  \int_0^{\pi} \sin^5\theta \, \md\theta = \frac{16}{15}
              \end{equation}
              所以
              \begin{equation}
                  \oiint_S x^4 \, \md S = 64 \cdot \frac{3\pi}{4} \cdot \frac{16}{15} = \frac{256\pi}{5}
              \end{equation}

              代回：
              \begin{equation}
                  \oiint_S \bm{F} \cdot \md\bm{S} = \frac12 \left[ 4 \cdot \frac{64\pi}{3} - \frac{256\pi}{5} \right]
                  = \frac12 \left[ \frac{256\pi}{3} - \frac{256\pi}{5} \right]
              \end{equation}
              \begin{equation}
                  = \frac12 \cdot 256\pi \cdot \frac{2}{15} = \frac{256\pi}{15}
              \end{equation}
              然后计算
              \begin{equation}
                  \iiint_V (\nabla \cdot \bm{F}) \, \md V
              \end{equation}
              计算散度符咒作用结果：
              \begin{equation}
                  \nabla \cdot \bm{F} = \frac{\partial}{\partial x}(xy^2) + \frac{\partial}{\partial y}(z\sin y) + \frac{\partial}{\partial z}(x^2 z) = y^2 + z\cos y + x^2
              \end{equation}
              由于对称性，$\iiint_V z\cos y \, \md V = 0$。所以
              \begin{equation}
                  \iiint_V (\nabla \cdot \bm{F}) \, \md V = \iiint_V (y^2 + x^2) \, \md V
              \end{equation}
              由对称性，$\iiint_V x^2 \, \md V = \iiint_V y^2 \, \md V$，且
              \begin{equation}
                  \iiint_V (x^2 + y^2 + z^2) \, \md V = \iiint_V r^2 \, \md V
              \end{equation}
              其中$r^2 = x^2 + y^2 + z^2$。在球坐标下，
              \begin{equation}
                  \iiint_V r^2 \, \md V = \int_0^{2\pi} \int_0^{\pi} \int_0^2 r^4 \sin\theta \, \md r \, \md\theta \, \md\phi
              \end{equation}
              \begin{equation}
                  = \int_0^{2\pi} \md\phi \int_0^{\pi} \sin\theta \, \md\theta \int_0^2 r^4 \, \md r = 2\pi \cdot 2 \cdot \frac{32}{5} = \frac{128\pi}{5}
              \end{equation}
              所以
              \begin{equation}
                  3 \iiint_V x^2 \, \md V = \frac{128\pi}{5} \quad\Rightarrow\quad \iiint_V x^2 \, \md V = \frac{128\pi}{15}
              \end{equation}
              因此
              \begin{equation}
                  \iiint_V (\nabla \cdot \mathbf{F}) \, \md V = 2 \cdot \frac{128\pi}{15} = \frac{256\pi}{15}
              \end{equation}
              可以看出：
              \begin{equation}
                  \oiint_{S}\bm{F}\cdot\md\bm{S} = \iiint_{V}(\nabla\cdot\bm{F})\,\md V = \frac{256\pi}{15}
              \end{equation}
              由此验证了高斯符咒公式把「曲面作用域」升格为「体积作用域」。

        \item 依据题意和斯托克斯符咒有
              \begin{equation}
                  \oint_{C}\bm{F}\cdot\md\bm{r}=\iint_{S}(\nabla\times\bm{F})\cdot\md\bm{S}。
              \end{equation}
              上式左侧为绫能直接作用的，右侧为绫想要作用的。\\
              先求旋度
              \begin{equation}
                  \nabla\times\bm{F}=(-\sin y,\,-2xz,\,-2xy)。
              \end{equation}
              故有：
              \begin{align}
                  \iint_{S}(\nabla\times\bm{F})\cdot\md\bm{S}
                   & = \iint_{x^{2}+y^{2}\leq 9}(-\sin y,\,-2x,\,-2xy)\cdot(0,0,1)\,\md x\md y \\
                   & = \iint_{x^{2}+y^{2}\leq 9} -2xy\,\md x\md y                              \\
                   & = 0
              \end{align}
              则作用的结果为$0$。\\
    \end{enumerate}
\end{example}

\section{世界理论}\label{sec:worldtheory}
本节用离散数学与抽象代数语言，给出“世界”的形式定义，并讨论其性质与层级结构。随后，基于此定义，分析“神”与“无全域神定理”等哲学命题在该世界观下的数学表达。
\begin{introduction}
    \item \hyperref[subsec:worldandproperty]{世界与世界的性质}
    \item \hyperref[subsec:levelofworld]{世界的层级}
    \item \hyperref[subsec:subworld]{结构子世界}
    \item \hyperref[subsec:god]{神}
    \item \hyperref[subsec:no-god]{无全域神定理}
    \item \hyperref[subsec:character-god]{角色神}
\end{introduction}

\subsection{世界与世界的性质}\label{subsec:worldandproperty}
在\wkn{}的世界观中，一个“\textbf{世界（World）}”是一个包含个体与命题的数学结构。个体是“存在的东西”，命题是“关于这些东西的陈述”。每个命题在某个世界中要么为真，要么为假。不同世界可以有不同的个体与命题集合。
基于此，给出如下世界的定义：

\begin{definition}[世界]\label{def:world}
    一个\textbf{世界（World）} $w$ 是三元组
    \begin{equation}
        w=(D_w,\Phi_w,\models_w)
    \end{equation}
    其中
    \begin{enumerate}
        \item $D_w\neq\varnothing$ 是\textbf{个体集合（Domain）}，称为世界的\textbf{域（Domain）}；
        \item $\Phi_w$ 是世界的\textbf{命题集合（Proposition Set）}，满足
              \begin{itemize}
                  \item 对任意 $\varphi\in\Phi_w$，其\textbf{否定（Negation）} $\neg\varphi\in\Phi_w$；
                  \item 对任意 $\varphi\in\Phi_w$，\textbf{排中律（Law of Excluded Middle）}成立：$\varphi\not\equiv\neg\varphi$；
              \end{itemize}
        \item $\models_w : \Phi_w \to \{\text{true}, \text{false}\}$ 是\textbf{真值指派函数（Truth Assignment Function）}，满足
              \begin{itemize}
                  \item 对任意 $\varphi\in\Phi_w$，$\models_w(\varphi)$ 仅取值于 \text{true} 或 \text{false}；
                  \item 对任意 $\varphi\in\Phi_w$，$\models_w(\neg\varphi)$ 与 $\models_w(\varphi)$ 互为否定；
                  \item $\models_w$ 在 $\Phi_w$ 上是唯一的。
              \end{itemize}
    \end{enumerate}
\end{definition}

\begin{remark}
    定义~\ref{def:world} 中的条件保证了每个世界的逻辑自洽性。个体集合 $D_w$ 保证了世界中“存在”的东西不为空；命题集 $\Phi_w$ 的闭包性质（对否定封闭）和排中律保证了逻辑的完备性与非矛盾性；真值指派函数 $\models_w$ 的唯一性确保了每个命题在该世界中的真值是确定的。
\end{remark}

\begin{remark}
    这样的定义并非完美。例如，该定义忽略了非命题性陈述（如对“存在某个东西”的“东西”的感知体验）以及模糊命题（如“这个东西很大”）。在命题的“本体论地位”问题上，该定义也难以给出很好的解释。
    确切的说，这个定义更像是在\wkn{}的世界观下，给出的一种理想化的“数学世界”模型。它适用于分析逻辑与集合论等基础数学问题，但在处理更复杂的哲学或语义问题时，可能需要更丰富的结构。
\end{remark}

\begin{example}[最小世界（Minimal World）]
    令 $D_w=\{a\}$，$\Phi_w=\{p,\neg p\}$，定义真值指派 $\models_w(p)=\text{true}$，则 $\models_w(\neg p)=\text{false}$。
    此时 $w$ 是一个合法世界，且 $|\Phi_w|=2$，$|D_w|=1$。
\end{example}

\begin{definition}[世界类]\label{def:W}
    令
    \begin{equation}
        \mathcal{W}=\{w\mid w\text{ 是满足定义~\ref{def:world} 下的世界}\}.
    \end{equation}
    若不限制世界中可使用的个体与命题来源，则在 ZFC 下，$\mathcal{W}$ 应理解为\textbf{真类（Proper Class）}，即不属于任何集合。这样的 $\mathcal{W}$ 包含所有可能的世界，称作\textbf{世界类（World Class）}。

    为了在后文中使用普通图论、偏序集与函数等集合论工具，我们固定一个足够大的\textbf{工作世界域（Working World Universe）} $\mathcal{W}_0\subseteq\mathcal{W}$，并要求 $\mathcal{W}_0$ 是集合。后文若无特别说明，所有世界图、层级关系、神的掌控范围等对象均在该固定集合 $\mathcal{W}_0$ 内讨论。
\end{definition}

\begin{theorem}[有限世界集合的基数上界]
    若固定一个含 $N$ 个候选个体的集合与一个含 $M$ 个候选命题的集合，并只考虑由它们生成的世界，则
    \begin{equation}
        \mathcal{W}_{N,M}=\{w\in\mathcal{W}_0\mid |D_w|\leq N,\,|\Phi_w|\leq M\}
    \end{equation}
    是有限集，其基数的一个上界为：
    \begin{equation}
        |\mathcal{W}_{N,M}| \leq \sum_{d=1}^{N} \sum_{p=0}^{M} \binom{N}{d} \cdot \binom{M}{p} \cdot 2^{p}.
    \end{equation}
    其中\(\dsy\binom{M}{p}\)表示从\(M\)个命题中选择\(p\)个命题的组合数。
\end{theorem}
\begin{proof}
    对于每个可能的个体数 \(d\)（从 1 到 \(N\)），以及每个可能的命题数 \(p\)（从 0 到 \(M\)），
    \begin{itemize}
        \item 个体集合 \(D_w\) 的选择有 \(\dsy\binom{N}{d}\) 种方式；
        \item 命题集合 \(\Phi_w\) 的选择有 \(\dsy\binom{M}{p}\) 种方式；
        \item 真值指派函数 \(\models_w\) 对于每个命题有 2 种选择（真或假），因此对于 \(p\) 个命题，有 \(2^p\) 种不同的指派方式。
    \end{itemize}
    因此，对于固定的 \(d\) 和 \(p\)，可能的世界数为：
    \begin{equation}
        \binom{N}{d} \cdot \binom{M}{p} \cdot 2^p.
    \end{equation}
    将所有可能的 \(d\) 和 \(p\) 的组合相加，得到总的世界数上界：
    \begin{equation}
        |\mathcal{W}_{N,M}| \leq \sum_{d=1}^{N} \sum_{p=0}^{M} \binom{N}{d} \cdot \binom{M}{p} \cdot 2^{p}.
    \end{equation}
    此上界是宽松的，因为在实际的合法世界中：
    \begin{enumerate}
        \item 命题集 \(\Phi_w\) 必须对否定封闭（即若 \(\varphi \in \Phi_w\)，则 \(\neg\varphi \in \Phi_w\)），这意味着 \(p\) 必须为偶数；
        \item 排中律要求命题与其否定不能等价，这进一步限制了真值指派的可能；
        \item 命题的内容应与个体域相关，而此处我们假设所有 \(M\) 个命题都是独立的。
    \end{enumerate}
    由于我们计算了所有可能的结构（包括那些不满足定义~\ref{def:world}中全部条件的结构），因此得到的是一个上界。
\end{proof}

\begin{theorem}[世界不可混淆性（World Non-Confusability）]\label{thm:distinct}
    对任意 $w_1,w_2\in\mathcal{W}_0$，若存在 $\varphi\in\Phi_{w_1}\cap\Phi_{w_2}$ 使得
    $\models_{w_1}(\varphi)\neq\;\models_{w_2}(\varphi)$，则 $w_1\neq w_2$。
\end{theorem}
\begin{proof}
    由定义~\ref{def:world}，一个世界的真值指派函数是唯一的。若 $w_1 = w_2$，则 $\models_{w_1}$ 和 $\models_{w_2}$ 是同一个函数，在 $\Phi_{w_1}\cap\Phi_{w_2}$ 上取值必然相同，与假设矛盾。
\end{proof}

\subsection{世界的层级}\label{subsec:levelofworld}
为了描述“一个世界比另一个世界拥有更多信息”这一直观想法，我们借用图论中的 \textbf{DAG（Directed Acyclic Graph， 有向无环图）}来描述世界之间的关系。DAG 的顶点是世界，边代表“命题扩张”的关系。

\begin{definition}[世界图与世界的层级关系]\label{def:graph}
    称二元组 $G=(\mathcal{W}_0,E)$ 为\textbf{世界图（World Graph）}，其中边集
    \begin{equation}
        E=\{(w,w')\in\mathcal{W}_0\times\mathcal{W}_0\mid \Phi_w \subsetneq \Phi_{w'},\ \text{且 } \models_{w'} \text{ 在 } \Phi_w \text{ 上的限制}\models_{w'}|_{\Phi_w}\text{等于 } \models_w\}.
    \end{equation}
    该图是严格的DAG，因而诱导出偏序集 $(\mathcal{W}_0,\preceq)$，其中
    \begin{equation}
        w\preceq w'\;\stackrel{\text{def}}{\Longleftrightarrow}\;
        \Phi_w \subseteq \Phi_{w'}\ \text{且 } \models_{w'} \text{ 在 } \Phi_w \text{ 上的限制}\models_{w'}|_{\Phi_w}\text{等于 } \models_w.
    \end{equation}
    此时称 $w$ 为 $w'$ 的\textbf{信息子世界（Information Subworld）}。
    若 $w \preceq w'$ 且 $w \neq w'$，则记作 $w \prec w'$，称 $w$ 为 $w'$ 的\textbf{真信息子世界（Proper Information Subworld）}。
    称偏序关系 $\preceq$ 为\textbf{世界的层级关系（Hierarchy Relation of Worlds）}。
\end{definition}

\begin{remark}
    函数的限制：若 $f:A\to B$ 是函数，$A'\subseteq A$，则 $f|_{A'}:A'\to B$ 是一个函数，满足对任意 $a\in A'$，有 $f|_{A'}(a)=f(a)$。
\end{remark}

\begin{remark}
    偏序关系：满足以下三条性质的二元关系，一般用$\preceq$表示：
    \begin{itemize}
        \item 自反性（Reflexivity）：对任意 $w\in\mathcal{W}_0$，有 $w\preceq w$；
        \item 反对称性（Antisymmetry）：对任意 $w_1,w_2\in\mathcal{W}_0$，若 $w_1\preceq w_2$ 且 $w_2\preceq w_1$，则 $w_1=w_2$；
        \item 传递性（Transitivity）：对任意 $w_1,w_2,w_3\in\mathcal{W}_0$，若 $w_1\preceq w_2$ 且 $w_2\preceq w_3$，则 $w_1\preceq w_3$。
    \end{itemize}
\end{remark}

\begin{remark}
    定义~\ref{def:graph} 中的边 $(w,w')$ 表示 $w'$ 在 $w$ 的基础上增加了新的命题，且对 $w$ 中已有命题的真值保持一致。偏序关系 $\preceq$ 则表示层级关系。
    由于命题集严格递增，世界图中不存在自环，从而保证了严格关系 $\prec$ 的反自反性，并避免在该层级关系中出现循环。
\end{remark}

\begin{definition}[世界的哈斯图]\label{def:hasse}
    世界 $G=(\mathcal{W}_0,E)$ 的\textbf{哈斯图（Hasse Diagram）}是一个简化的有向图，去掉了所有可由传递性推导出的边。
    即若存在 $w_1,w_2,w_3\in\mathcal{W}_0$ 使得 $w_1 \prec w_2$ 且 $w_2 \prec w_3$，则在哈斯图中不画出边 $(w_1,w_3)$，
    并省略箭头。若 $w_1 \prec w_2$，则在哈斯图中画出一条从 $w_1$ 指向 $w_2$ 的线段，且 $w_2$ 位于 $w_1$ 的\textbf{上方}。
\end{definition}
\begin{example}
    下面是一个世界图与世界的哈斯图示例。设 $w_0,w_1,w_2,w_3$ 为四个世界，满足：
    \begin{equation}
        w_0 \prec w_1,\quad w_0 \prec w_2,\quad w_1 \prec w_3,\quad w_2 \prec w_3.
    \end{equation}
    \begin{center}
        \begin{tikzpicture}[node distance=1.8cm, auto]
            \node (w0) [circle, draw] {$w_0$};
            \node (w1) [circle, draw, right=of w0] {$w_1$};
            \node (w2) [circle, draw, below=of w0] {$w_2$};
            \node (w3) [circle, draw, below=of w1] {$w_3$};
            \draw[->] (w0) to (w1);
            \draw[->] (w0) to (w2);
            \draw[->] (w1) to (w3);
            \draw[->] (w2) to (w3);
            \draw[->] (w0) to (w3);
        \end{tikzpicture}
        \qquad
        \begin{tikzpicture}[node distance=1.8cm, auto]
            \node (v3) [circle, draw] {$w_3$};
            \node (v1) [circle, draw, below left=1.2cm and 0.8cm of v3] {$w_1$};
            \node (v2) [circle, draw, below right=1.2cm and 0.8cm of v3] {$w_2$};
            \node (v0) [circle, draw, below=2.4cm of v3] {$w_0$};
            \draw (v0) to (v1);
            \draw (v0) to (v2);
            \draw (v1) to (v3);
            \draw (v2) to (v3);
        \end{tikzpicture}
    \end{center}
    可以看到，左图为世界图，右图为哈斯图，二者的区别在于去掉了由偏序关系的传递性推导出的边 $(w_0,w_3)$，
    并去掉了箭头。此外，哈斯图中更强调了层级关系，$w_3$ 位于最高层，$w_0$ 位于最低层，对应于信息子世界的层级关系。
\end{example}

\begin{lemma}[良基子集的极小元存在]\label{lemma:min}
    设 $S\subseteq\mathcal{W}_0$ 为非空子集。若 $S$ 在关系 $\prec$ 下不存在无限严格下降链
    \begin{equation}
        \cdots \prec w_2 \prec w_1 \prec w_0,
    \end{equation}
    则 $S$ 中存在 $\preceq$-极小元。
\end{lemma}
\begin{proof}
    假设 $S$ 没有极小元。任取 $w_0\in S$，由于 $w_0$ 不是极小元，存在 $w_1\in S$ 使得 $w_1\prec w_0$。同理可递归构造 $w_{n+1}\prec w_n$，从而得到无限严格下降链
    \begin{equation}
        \cdots \prec w_2 \prec w_1 \prec w_0,
    \end{equation}
    这与假设矛盾。故 $S$ 必有极小元。
\end{proof}

\begin{remark}
    若不加入良基条件，上述结论一般不成立。例如命题集可以形成无限真包含下降链，此时相应的世界子集可能没有极小元。因此，层级函数只在满足良基条件的世界域或分支上具有通常意义。
\end{remark}

\begin{remark}
    \textbf{极小元（Minimal Element）}是指在某子集中没有比它更“信息更少”的世界。注意，与最小元相比，极小元不一定唯一，且可能有多个极小元位于不同的连通分量上。

    例如，若 $w_1$ 和 $w_2$ 分别是两个不相关分支的极小元，则它们都没有比自己信息更少的世界，都是极小元，但 $w_1 \not\preceq w_2$ 且 $w_2 \not\preceq w_1$。

    遗憾的是，世界一般不存在\textbf{最小元（Least Element）}。因为世界之间的关系是偏序关系而非全序关系，任意两个世界不一定可比。
\end{remark}

\begin{definition}[层函数]\label{def:level}
    在满足引理~\ref{lemma:min} 的良基工作世界域中，定义\textbf{层函数（Level Function）} $\mathrm{level}:\mathcal{W}_0\to\mathbb{N}\cup\{\infty\}$ 如下：
    \begin{equation}
        \mathrm{level}(w)=
        \begin{cases}
            \min \big\{ n \in \mathbb{N} \mid \exists w_0 \preceq w_1 \preceq \dots \preceq w_n = w,\ w_0 \text{ 是 } \mathcal{W}_0 \text{ 的极小元} \big\}, & \text{如果存在这样的链}, \\
            \infty,                                                                                                                                   & \text{否则}.
        \end{cases}
    \end{equation}
    即 $\mathrm{level}(w)$ 是从某个极小元到 $w$ 的最短链长度；若 $w$ 不可从任何极小元到达，则层数为 $\infty$。
\end{definition}

\begin{remark}
    层函数把世界图划分成“代际”，但允许不同连通分量的世界拥有任意有限的层数或 $\infty$。
    第 0 层由所有极小元组成（可能有多个，且位于不同连通分量）。
    若 $w \prec w'$，则 $\mathrm{level}(w) < \mathrm{level}(w')$；但反之，层数小的世界不一定可比（例如不同分支的世界层数可能相同但毫不相关）。
\end{remark}

\begin{example}\label{ex:level}
    考虑三个世界 $w_0, w_1, w_2$，满足：
    \begin{itemize}
        \item $w_0$：$D=\{a\}$，$\Phi_{w_0} = \{P(a), \neg P(a)\}$，$\models_{w_0}(P(a))=\text{true}$。
        \item $w_1$：在 $w_0$ 上增加命题 $Q(a)$，且 $\models_{w_1}(Q(a))=\text{false}$。
        \item $w_2$：$D=\{b\}$，$\Phi_{w_2} = \{R(b), \neg R(b)\}$，$\models_{w_2}(R(b))=\text{true}$（与 $w_0, w_1$ 无关）。
    \end{itemize}
    则 $w_0 \prec w_1$，且 $w_0, w_2$ 都是极小元。
    $\mathrm{level}(w_0)=0$，$\mathrm{level}(w_1)=1$，$\mathrm{level}(w_2)=0$。
    $w_1$ 与 $w_2$ 层数不同且不可比，体现了不同分支的世界可以“层级毫不相关”。
\end{example}

\subsection{结构子世界}\label{subsec:subworld}
除了信息子世界，我们还可以定义“\textbf{结构子世界（Structural Subworld）}”这一更强的概念。
在直观上，一个“结构子世界”应当是大世界的一部分——不仅拥有的信息更少，其个体域也应包含于大世界的个体域中，且关于这部分个体的所有真值均与大世界一致。

\begin{definition}[结构子世界]\label{def:subworld}
    设 $w = (D_w, \Phi_w, \models_w)$ 和 $w' = (D_{w'}, \Phi_{w'}, \models_{w'})$ 是两个世界。
    称 $w$ 是 $w'$ 的\textbf{结构子世界（Structural Subworld）}，记作 $w \sqsubseteq w'$，当且仅当以下条件同时成立：
    \begin{enumerate}
        \item \textbf{个体域包含（Domain Inclusion）}：$D_w \subseteq D_{w'}$。
        \item \textbf{命题集一致性（Proposition Set Consistency）}：$\Phi_w = \{\, \varphi \in \Phi_{w'} \mid \text{$\varphi$ 中出现的所有个体均属于 } D_w \,\}$。
        \item \textbf{真值继承性（Truth Inheritance）}：对任意 $\varphi \in \Phi_w$，有 $\models_w(\varphi) = \models_{w'}(\varphi)$。
    \end{enumerate}
    若 $w \sqsubseteq w'$ 且 $w \neq w'$，则记作 $w \sqsubset w'$，并称 $w$ 是 $w'$ 的\textbf{真结构子世界（Proper Structural Subworld）}。
\end{definition}


\begin{remark}
    定义~\ref{def:subworld} 中的条件 2 确保了 $w$ 的命题集恰好是 $w'$ 的命题集中那些“只谈论 $D_w$ 中个体”的命题。条件 3 则保证了对这些命题的真值判断，结构子世界与父世界完全一致。这表明，结构子世界完全由父世界通过限制个体域而得到。
\end{remark}

\begin{example}\label{ex:subworld}
    考虑世界 $w'$，其中 $D_{w'} = \{a, b, c\}$，命题集 $\Phi_{w'}$ 包含原子命题 $P(a), Q(b), R(a,c)$ 及其否定。真值指派 $\models_{w'}$ 满足 $\models_{w'}(P(a)) = \text{true}$，$\models_{w'}(Q(b)) = \text{false}$，$\models_{w'}(R(a,c)) = \text{true}$。

    现取 $D_w = \{a, c\}$。根据定义~\ref{def:subworld}，构造 $w$ 如下：
    \begin{itemize}
        \item $\Phi_w = \{\, \varphi \in \Phi_{w'} \mid \text{$\varphi$ 中只出现 $a, c$} \,\} = \{P(a), \neg P(a), R(a,c), \neg R(a,c)\}$。
        \item $\models_w(P(a)) = \models_{w'}(P(a)) = \text{true}$，$\models_w(R(a,c)) = \models_{w'}(R(a,c)) = \text{true}$，并相应地对否定命题赋值。
    \end{itemize}
    则 $w$ 是 $w'$ 的一个结构子世界，即 $w \sqsubseteq w'$。
\end{example}

\begin{theorem}[结构子世界与信息子世界的关系（Relation between Structural and Information Subworld）]\label{lemma:sub-info}
    若 $w \sqsubseteq w'$，则 $w \preceq w'$。即结构子世界必是信息子世界。
\end{theorem}
\begin{proof}
    由定义~\ref{def:subworld}，$D_w \subseteq D_{w'}$ 且 $\Phi_w \subseteq \Phi_{w'}$（因为 $\Phi_w$ 是 $\Phi_{w'}$ 的子集），并且真值在 $\Phi_w$ 上一致。这正好满足定义~\ref{def:graph} 中 $w \preceq w'$ 的条件。
\end{proof}

\begin{remark}
    反之不成立。信息子世界只要求命题集的包含与真值一致，但允许个体域完全不同。因此，结构子世界是信息子世界的一个特例，它强化了“部分结构”的直观。
\end{remark}

此外，我们将信息子世界与结构子世界统称为\textbf{子世界（Subworld）}。

\subsection{神}\label{subsec:god}
在许多哲学与神学体系中，“神”被描述为一个全知全能的存在，能够掌控一切事物。
基于世界与子世界的框架，我们可以对“神”这一概念进行数学建模。神被定义为一个能“掌控”某个主世界及其所有结构子世界的智能体。这种“掌控”具体体现为对其掌控范围内所有事实的无误且完全的认知。

\begin{definition}[神]\label{def:god}
    \textbf{神（God）} \(G\) 是一个四元组：
    \begin{equation}
        G = (w_G, \, \mathcal{S}_G, \, \pi, \, \Vdash_G)
    \end{equation}
    其中：
    \begin{enumerate}
        \item \( w_G = (D_{w_G}, \Phi_{w_G}, \models_{w_G}) \in \mathcal{W}_0 \) 是一个世界，称为\textbf{神 \(G\) 的主世界（Main World of God）}。
        \item \( \mathcal{S}_G = \{ w' \in \mathcal{W}_0 \mid D_{w'} \subseteq D_{w_G},\ \Phi_{w'} = \{\varphi \in \Phi_{w_G} \mid \text{$\varphi$中所有个体属于 } D_{w'}\},\ \models_{w'} = \models_{w_G}|_{\Phi_{w'}} \} \)
              是 \( w_G \) 的所有结构子世界的集合，称为\textbf{神 \(G\) 的掌控世界集（Control World Set of God \(G\)）}。
        \item \( \pi: \mathcal{S}_G \to \mathcal{P}(\Phi_{w_G}) \) 是一个\textbf{投影函数（Projection Function）}，满足 \( \pi(w') = \Phi_{w'} \)。
        \item \( \Vdash_G \subseteq \{ (w', \varphi) \in \mathcal{S}_G \times \Phi_{w_G} \mid \varphi \in \pi(w') \} \) 是一个\textbf{全局认知关系（Global Cognition Relation）}，满足：
              \begin{itemize}
                  \item \textbf{一致性公理（Consistency Axiom）}：若 \( (w', \varphi) \in \Vdash_G \)，则 \( \models_{w'}(\varphi) = \text{true} \)。
                  \item \textbf{全知公理（Omniscience Axiom）}：对任意 \( w' \in \mathcal{S}_G \) 和任意 \( \varphi \in \pi(w') \)，有 \( (w', \varphi) \in \Vdash_G \) 当且仅当 \( \models_{w'}(\varphi) = \text{true} \)。
              \end{itemize}
    \end{enumerate}
\end{definition}

\begin{remark}
    这个定义的核心在于将“掌控”具体化为两种能力：
    \begin{enumerate}
        \item \textbf{范围}：神的掌控范围是其主世界的所有结构子世界 \( \mathcal{S}_G \)。
        \item \textbf{认知}：神拥有一个全局的、无误的认知关系 \( \Vdash_G \)，能够无误地知晓其掌控范围内任何局部（任何结构子世界）的任何事实。
    \end{enumerate}
\end{remark}

\begin{remark}
    区分宗教意义上的“神”与\wkn 世界观中的“神”是重要的。这里的“神”是一个理想化的数学构造，旨在捕捉全知全能的概念，而非特定宗教或文化中的神学定义。
\end{remark}

基于这个定义，我们可以推导出神的一些重要性质。

\begin{theorem}[神对主世界的全知]
    神对其主世界 \( w_G \) 是全知的。即，对于任意 \( \varphi \in \Phi_{w_G} \)，有：
    \begin{equation}
        (w_G, \varphi) \in \Vdash_G \quad \text{当且仅当} \quad \models_{w_G} (\varphi) = \text{true}.
    \end{equation}
\end{theorem}

\begin{theorem}[神对局部信息的掌控]
    设 \( w' \in \mathcal{S}_G \) 且 \( w'' \in \{w | w \in \mathcal{W}_0 \land w \sqsubseteq w'\}\)，则 \( w'' \in \mathcal{S}_G \)。进而，神 \(G\) 对 \( w'' \) 中的任意命题 \( \varphi \) 的认知 \( (w'', \varphi) \in \Vdash_G \) 与 \( w' \) 中对 \( \varphi \) 的认知 \( (w', \varphi) \in \Vdash_G \) 是一致的（当 \( \varphi \in \pi(w'') \) 时），且都符合各自世界的客观真值。

\end{theorem}
\begin{proof}
    由子世界关系的传递性可得 \( w'' \in \mathcal{S}_G \)。其认知的一致性由全知公理保证，因为 \( \models_{w''}(\varphi) \) 与 \( \models_{w'}(\varphi) \) 在 \( \varphi \in \pi(w'') \) 时是一致的（根据结构子世界的定义）。
\end{proof}

\begin{example}
    假设主世界 \( w_G \)满足：
    \begin{equation} D_{w_G} = \{a, b\} \end{equation}
    \begin{equation} \Phi_{w_G} = \{P(a), \neg P(a), Q(b), \neg Q(b), R(a,b), \neg R(a,b)\} \end{equation}
    其真值指派为
    \begin{equation}
        \models_{w_G}(P(a)) = \text{true},\; \models_{w_G}(Q(b)) = \text{false},\; \models_{w_G}(R(a,b)) = \text{true}.
    \end{equation}

    那么：
    \begin{itemize}
        \item \( \mathcal{S}_G \) 将包含：
              \begin{itemize}
                  \item \( w_G \) 自身。
                  \item 个体域为 \( \{a\} \) 的结构子世界 \( w_a \)，其中 \( \Phi_{w_a} = \{P(a), \neg P(a)\} \)，\( \models_{w_a}(P(a)) = \text{true} \)。
                  \item 个体域为 \( \{b\} \) 的结构子世界 \( w_b \)，其中 \( \Phi_{w_b} = \{Q(b), \neg Q(b)\} \)，\( \models_{w_b}(Q(b)) = \text{false} \)。
              \end{itemize}
        \item 函数 \( \pi \) 满足 \( \pi(w_G) = \Phi_{w_G} \), \( \pi(w_a) = \Phi_{w_a} \), \( \pi(w_b) = \Phi_{w_b} \)。
        \item 关系 \( \Vdash_G \) 则包含了诸如：
              \begin{itemize}
                  \item \( (w_G, P(a)) \in \Vdash_G \), \( (w_G, R(a,b)) \in \Vdash_G \)（因为它们在 \( w_G \) 中为真）。
                  \item \( (w_a, P(a)) \in \Vdash_G \)（因为它在 \( w_a \) 中为真）。
                  \item \( (w_b, \neg Q(b)) \in \Vdash_G \)（因为 \( Q(b) \) 在 \( w_b \) 中为假）。
                  \item \( (w_a, R(a,b)) \notin \Vdash_G \)（因为 \( R(a,b) \notin \pi(w_a) \)）。
              \end{itemize}
    \end{itemize}
    神 \( G \) 因此同时知晓全局状态 \( w_G \) 和所有局部状态 \( w_a, w_b \) 的全部真理。
\end{example}

\subsection{无全域神定理}\label{subsec:no-god}

基于已建立的世界层级结构与神的定义，我们可以提出一个更为根本的论断：在命题可以无限扩张的假设下，掌控所有世界的神在逻辑上是不可能的。此结论不依赖于个体域或具体的真值冲突，仅源于世界偏序集本身的结构性质。

\begin{axiom}[命题的无限扩张性]
    工作世界域 \( \mathcal{W}_0 \) 中的命题集可以无限扩张，这就是\textbf{命题的无限扩张性公理（Infinite Expandability of Propositions Axiom）}。形式化地：
    \begin{equation}
        \forall w \in \mathcal{W}_0,\; \exists w' \in \mathcal{W}_0,\; \text{使得 } \Phi_w \subsetneq \Phi_{w'}。
    \end{equation}
    该公理断言，不存在一个拥有“所有”命题的世界。
\end{axiom}

\begin{definition}[全域神]
    设 \( G^* = (w_{G^*}, \mathcal{S}_{G^*}, \pi, \Vdash_{G^*}) \) 是一个神。若 \( \mathcal{S}_{G^*} = \mathcal{W}_0 \)，即其掌控范围包含工作世界域中的所有世界，则称 \( G^* \) 为\textbf{全域神（Universal God）}。
\end{definition}


\begin{theorem}[无全域神定理（No Universal God Theorem）]\label{thm:no-god}
    在 ZFC、固定工作世界域与命题的无限扩张性公理下，不存在全域神。
\end{theorem}

\begin{proof}
    采用反证法。
    \begin{enumerate}
        \item 假设存在一个全域神 \( G^* = (w_{G^*}, \mathcal{S}_{G^*}, \pi, \Vdash_{G^*}) \)。根据全域神的定义，有 \( \mathcal{S}_{G^*} = \mathcal{W}_0 \)。
        \item 由于 \( \mathcal{S}_{G^*} \) 是 \( w_{G^*} \) 的所有结构子世界组成的集合，而 \( \mathcal{S}_{G^*}=\mathcal{W}_0 \)，这意味着：
              \begin{equation}
                  \forall w \in \mathcal{W}_0,\; D_w \subseteq D_{w_{G^*}} \ \text{且} \ \Phi_w \subseteq \Phi_{w_{G^*}}。
              \end{equation}
              特别地，\( w_{G^*} \) 的命题集 \( \Phi_{w_{G^*}} \) 必须包含所有其他世界的命题集。
        \item 现在，考虑神 \( G^* \) 的主世界 \( w_{G^*} \)。根据\textbf{命题的无限扩张性公理（Infinite Expandability of Propositions）}，存在另一个世界 \( w' \in \mathcal{W}_0 \)，使得：
              \begin{equation}
                  \Phi_{w_{G^*}} \subsetneq \Phi_{w'}。
              \end{equation}
        \item 但是根据第2点，由于 \( w' \in \mathcal{W}_0 \)，必须有 \( \Phi_{w'} \subseteq \Phi_{w_{G^*}} \)。
        \item 这与第3点得出的 \( \Phi_{w_{G^*}} \subsetneq \Phi_{w'} \) 相矛盾，因为一个集合不能既是另一个集合的真超集，又是其子集。
        \item 因此，最初的假设（存在全域神）是错误的。故全域神 \( G^* \) 不存在。
    \end{enumerate}
\end{proof}

\begin{remark}
    此证明的不关心世界的具体内容（个体、真值），仅利用了两个基本事实：
    \begin{enumerate}
        \item 全域神的存在蕴含了存在一个世界，其命题集包含所有其他世界的命题集。
        \item 命题的无限扩张性保证了这样的世界不可能存在。
    \end{enumerate}
    全域神因其试图以某个主世界为基底掌控整个工作世界域，而该工作世界域在命题无限扩张假设下不存在命题最大的世界。
\end{remark}

\begin{corollary}[多神论的逻辑必然性]
    任何具有神的掌控范围必然是局部的。即，存在多个神 \( G_1, G_2, \dots, G_n \)，分别对应于不同的主世界 \( w_{G_1}, w_{G_2}, \dots, w_{G_n} \)，且满足：
    \begin{equation}
        \mathcal{S}_{G_i} \subsetneq \mathcal{W}_0, \quad \bigcup_{i=1}^n \mathcal{S}_{G_i} \subseteq \mathcal{W}_0。
    \end{equation}
    从逻辑上讲，\textbf{多神论（Polytheism）}是此框架下的必然推论。
\end{corollary}


无全域神定理并未断言“神”不存在，而是精确地证明了“全域神”这一概念的逻辑不自洽。它为我们描绘了一幅图景：神性是局部的、分层的，与特定的世界绑定。绝对的、统摄一切的唯一主宰，在逻辑的审视下，只是一个无法实现的幻象。


\subsection{角色神}\label{subsec:character-god}

在叙事作品中，作者常常塑造一些拥有近乎神明般力量的角色。它们可能在其所处的故事范围内无所不能，但作者又会通过剧情明确展示其能力的边界（例如，无法改变某种宇宙法则、无法干涉另一个维度、或者存在一个它无法解答的问题）。为了在形式框架内精确捕捉这一概念，并与我们之前定义的“神”相区分，我们引入“角色神（Character God）”的概念。

\begin{definition}[角色神（Character God）]\label{def:character-god}
    \textbf{角色神（Character God）}是神结构的叙事性弱化。它仍是一个四元组
    \begin{equation}
        G = (w_G, \mathcal{S}_G, \pi, \Vdash_G),
    \end{equation}
    其中 \(w_G\)、\(\mathcal{S}_G\)、\(\pi\) 的定义与定义~\ref{def:god} 相同，但认知关系 \(\Vdash_G\) 不再要求满足全知公理，而只要求满足一致性公理：若 \((w',\varphi)\in\Vdash_G\)，则 \(\models_{w'}(\varphi)=\text{true}\)。

    若存在一个非空的\textbf{限制集（Restriction Set）} \( R_G \subseteq \Phi_{w_G} \)，使得其中每个真命题都不被角色神认知，即
    \begin{equation}
        \forall \varphi \in R_G,\; \models_{w_G}(\varphi)=\text{true}\Rightarrow (w_G, \varphi) \notin \Vdash_G,
    \end{equation}
    则称 \( G \) 为一个\textbf{角色神（Character God）}。
    集合 \( R_G \) 中的命题称为该角色神的\textbf{未知命题（Unknown Proposition）}或\textbf{限制命题（Restricted Proposition）}。
\end{definition}

\begin{remark}
    角色神的定义体现了对“神性（Divinity）”的一种\textbf{叙事性削弱}：
    \begin{itemize}
        \item 一个角色神在结构上与神相似：它拥有主世界、掌控范围、投影函数和认知关系。
        \item 它与神的关键区别在于全知公理被削弱。角色神对其主世界并非全知，它存在一个明确的认知盲区 \( R_G \)。
    \end{itemize}
\end{remark}

\begin{theorem}[角色神的性质]
    设 \( G \) 是一个角色神，其限制集为 \( R_G \)。
    \begin{enumerate}
        \item \( G \) 对其主世界 \textbf{不是全知的（Not Omniscient）}。即，\( \exists \varphi \in \Phi_{w_G} \) 使得 \( (w_G, \varphi) \notin \Vdash_G \)。
        \item \( G \) 的认知关系 \( \Vdash_G \) 在其掌控范围 \( \mathcal{S}_G \) 上仍然是\textbf{一致的（Consistent）}（不会认知错误），但是\textbf{不完全的（Incomplete）}（存在不知道的真命题）。
        \item 若 \( R_G = \Phi_{w_G} \)，则 \( G \) 对其主世界一无所知。这是一种极端情况，可称为\textbf{无知的神（Ignorant God）}。
    \end{enumerate}
\end{theorem}

\begin{definition}[神性强度]
    我们可以用限制集的大小来粗略衡量一个角色神的\textbf{“神性强度（Strength of Divinity）”}：
    \begin{itemize}
        \item 若 \( |R_G| = 0 \)，则 \( G \) 是真正的\textbf{神（God）}。
        \item 若 \( 0 < |R_G| \ll |\Phi_{w_G}| \)，则 \( G \) 是\textbf{近乎全能者（Nearly Omnipotent Being）}。
        \item 若 \( |R_G| \approx |\Phi_{w_G}| \)，则 \( G \) 是\textbf{极度受限者（Extremely Restricted Being）}。
    \end{itemize}
    这种度量方式为比较不同叙事中“神明角色”的能力提供了形式化的标尺。
\end{definition}

\begin{theorem}[无角色神定理]
    \textbf{全域角色神（Universal Character God）}（即满足 \( \mathcal{S}_G = \mathcal{W}_0 \) 的角色神）同样不存在。

\end{theorem}
\begin{proof}
    假设存在一个全域角色神 \( G^* \)。根据定义，有 \( \mathcal{S}_{G^*} = \mathcal{W}_0 \)。
    如定理~\ref{thm:no-god} 的证明，这要求 \( w_{G^*} \) 的命题集包含所有其他世界的命题集。
    但根据命题的无限扩张性公理，这样的世界不可能存在。
    因此，全域角色神 \( G^* \) 也不存在。
\end{proof}


“角色神”概念的引入，完成了我们对“神性”的形式化构建。
它允许我们在不违背无全域神定理的前提下，严谨地分析和讨论叙事作品中那些拥有巨大能力但却存在明确边界的“类神”角色。一个角色是否是“神”，不再是一个模糊的文学评论，而是一个可以通过分析其认知边界 \( R_G \) 来精确判定的数学问题。

\section{元素理论}\label{sec:elementtheory}
在\wkn{}世界中，元素学是解释物质性质、能量表现、生命倾向和术式来源的基础学说。这里的“元素”不是现实化学元素，而是世界内部用于描述万物倾向的分类体系。

元素学解释了许多常见现象：为什么某些地区多火山，为什么某些人擅长治疗术，为什么雷系术式常与速度和机械有关，为什么暗元素不一定代表邪恶，而常常代表遮蔽、吸收和边界。

\begin{introduction}
    \item \hyperref[subsec:elementorigin]{元素本源与存在论基础}
    \item \hyperref[subsec:elementproperties]{七种基础元素的属性谱系}
    \item \hyperref[subsec:elementinteraction]{元素交互的基本规则}
    \item \hyperref[subsec:elementinfluence]{元素理论对世界架构的影响}
\end{introduction}

\subsection{元素本源与存在论基础}\label{subsec:elementorigin}

\begin{remark}
    本节所述元素理论属于\wkn{}世界内部的\textbf{架空物理公理体系}。其中借用“场”“能量”“频率”“粒子”等现实科学术语作为类比语言，但这些术语的含义以本世界内部的元素学定义为准，不等同于现实物理学中的严格定义。
\end{remark}

元素学的基本观点并不是“世界由七种现实物质拼成”，而是：世界中一切物质、能量、生命和能力，都可以表现出某种元素倾向。一个石块不是只含土元素，一团火也不是只含火元素；它们只是分别表现出更强的“稳定倾向”和“释放倾向”。

为了描述这种倾向，元素学引入三个基础概念：
\begin{definition}[元子]
    \textbf{元子（Elementon）}是元素性质的最小单位。它不是现实物理中的基本粒子，而是元素学中用于解释元素现象的基本单位。若某个对象表现出明显的元素特征，就可以说它含有相应元素的元子结构。
\end{definition}

\begin{definition}[元素印记]
    \textbf{元素印记（Elemental Sigil）}是元子所携带的性质标识。它决定该元子更容易表现出稳定、流动、传递、释放、激活、显现或遮蔽等倾向。
\end{definition}

\begin{definition}[初始以太场]
    \textbf{初始以太场（Prima Aether Field）}是七种元素共同来源的背景场。它并不是某一种元素，而是元素能够存在、转化和相互影响的基础环境。
\end{definition}

用更直观的话说：元子是元素的“单位”，元素印记是元子的“性格”，初始以太场是元素活动的“舞台”。

\subsubsection{元素不是阵营}

元素只表示性质，不直接表示善恶。光元素不一定善良，暗元素也不一定邪恶。光可以代表显现、秩序和净化，也可以代表审判、暴露和压迫；暗可以代表隐藏、吸收和边界，也可以代表庇护、休眠和封印。

因此，元素亲和性不能直接等同于性格。火亲和者不一定暴躁，水亲和者不一定温柔，暗亲和者也不一定危险。元素提供的是倾向，而一个人如何表现这种倾向，取决于性格、环境、经历和术式训练。

\subsection{七种基础元素的属性谱系}\label{subsec:elementproperties}

七种基础元素分别是土、水、风、火、雷、光、暗。它们可以分为两组：
\begin{itemize}
    \item \textbf{实体倾向元素}：土、水、风。它们更接近世界中可感知、可承载、可流动的部分。
    \item \textbf{能量倾向元素}：火、雷、光、暗。它们更接近释放、激活、显现、抑制等变化过程。
\end{itemize}

这种分类不代表两组元素互不相通。水可以承载能量，火也可以改变物质结构；风可以传递声音，光也可以承载信息。

\begin{table}[htbp]
    \centering
    \caption{七种基础元素的基本定位}
    \label{tab:element-basic}
    \begin{tabular}{llll}
        \toprule
        \textbf{元素} & \textbf{名称} & \textbf{核心关键词} & \textbf{常见表现} \\
        \midrule
        土 & Tellus & 稳定、承载、结构 & 岩石、金属、骨骼、防御、城塞 \\
        水 & Aqua & 流动、修复、溶解 & 河流、血液、治疗、净化、适应 \\
        风 & Ventus & 传递、速度、扩散 & 气流、声音、移动、侦察、通信 \\
        火 & Ignis & 释放、燃烧、爆发 & 火焰、热量、锻造、破坏、激情 \\
        雷 & Fulgur & 激活、贯通、冲击 & 闪电、神经、机械、瞬发、麻痹 \\
        光 & Lux & 显现、秩序、净化 & 照明、感知、结界、审判、记录 \\
        暗 & Umbra & 遮蔽、吸收、边界 & 影子、封印、隐匿、沉眠、隔绝 \\
        \bottomrule
    \end{tabular}
\end{table}

\subsubsection{土元素：稳定与承载}

土元素代表结构、重量、支撑和持久。它常出现在山脉、矿物、城墙、骨骼、甲壳和大型防御术式中。土元素能力通常不追求速度，而追求可靠性。

土亲和者常承担防御、保护、建造、封锁等职责。其弱点是反应较慢，面对高机动敌人时容易被绕开。若土元素过强，个体可能变得迟钝、固执，或难以适应快速变化的局面。

\subsubsection{水元素：流动与修复}

水元素代表流动、适应、溶解和恢复。它常出现在河流、海洋、血液、药剂、治疗术和净化仪式中。水元素不一定柔弱，它也可以通过侵蚀、冲击和冻结形成强大的破坏力。

水亲和者常见于治疗、辅助、净化、潜入和环境控制等领域。其弱点是容易受容器和地形影响：没有水源、空间干燥或环境被污染时，水元素能力会变得不稳定。

\subsubsection{风元素：传递与速度}

风元素代表移动、扩散、声音和信息传递。它常出现在气流、飞行、远距通信、侦察术和速度强化中。风元素的重点不是“空气本身”，而是“从一处到另一处”的过程。

风亲和者常见于机动、侦察、支援和扰乱等领域。其弱点是难以长时间维持高强度实体防御，也容易被封闭空间、重力术式或土元素结构限制。

\subsubsection{火元素：释放与爆发}

火元素代表能量释放、燃烧、热量和变化。它常出现在火焰、爆炸、熔炉、锻造、战斗魔法和情绪爆发中。火元素常被用于攻击性术式，但它的意义不只是破坏，也包括点燃、驱寒、冶炼和重生。

火亲和者常见于爆发输出、范围攻击、锻造和精神鼓舞等领域。其弱点是消耗大、难以精细控制，且容易被水元素、缺氧环境或吸收型暗元素压制。

\subsubsection{雷元素：激活与贯通}

雷元素代表瞬间激活、贯通、冲击和连锁反应。它常出现在闪电、神经信号、机械启动、磁场、麻痹和高速打击中。雷元素与火元素都具有攻击性，但火偏持续释放，雷偏瞬间击穿。

雷亲和者常见于高速战斗、机械操作、神经强化和打断敌方术式等领域。其弱点是稳定性差，容易误伤，也容易受到绝缘、接地、干扰和精密控制不足的限制。

\subsubsection{光元素：显现与秩序}

光元素代表显现、辨识、净化、记录和秩序。它常出现在照明、观测、结界、神圣术式、记忆保存和真相揭示中。光元素并不等同于善，它也可能代表监视、审判和无法逃避的暴露。

光亲和者常见于侦测、净化、结界、治疗辅助和反隐能力等领域。其弱点是容易被遮蔽、折射、污染或暗元素吸收；过强的光也可能造成灼伤、失明或精神压迫。

\subsubsection{暗元素：遮蔽与边界}

暗元素代表隐藏、吸收、沉默、休眠和边界。它常出现在影子、封印、潜行、隔绝空间、诅咒、梦境和地下结构中。暗元素不是“坏元素”，它更像是让事物暂时离开视野、离开流动、离开作用范围的力量。

暗亲和者常见于潜入、封印、削弱、隔绝和情报保护等领域。其弱点是正面爆发通常不如火雷，且在强光、公开环境或秩序结界中容易被限制。

\subsection{元素交互的基本规则}\label{subsec:elementinteraction}

元素之间不是简单的固定克制表。更准确地说，元素交互取决于三个因素：元素本身的性质、环境条件、使用者的控制方式。同样是水与火，少量水可能被蒸发，大量水可以灭火；稳定的风可以助燃，过强的风也可以吹散火焰。

元素学通常将元素交互分为三类：相生、相克和转化。

\subsubsection{相生：帮助对方发挥}

\textbf{相生作用}指一种元素为另一种元素提供条件，使后者更容易发挥效果。相生并不等于两种元素完全相同，而是它们在某个目标上形成配合。

常见例子包括：
\begin{itemize}
    \item 风助火：稳定气流可以扩大火势，使火焰传播更快。
    \item 土蓄水：土元素结构可以形成河床、水库或地下水脉，使水元素更容易聚集。
    \item 光导雷：光元素的定位与显现能力可以帮助雷元素锁定目标。
    \item 暗养土：暗元素的沉眠与隔绝可以保护土元素结构，使封印、陵墓和地下城长期稳定。
\end{itemize}

\subsubsection{相克：限制对方发挥}

\textbf{相克作用}指一种元素削弱、阻断或扭曲另一种元素的效果。相克不是绝对胜负，而是“在当前条件下更容易压制”。

常见例子包括：
\begin{itemize}
    \item 水克火：水吸收热量并隔绝燃烧条件，使火元素难以持续。
    \item 土克风：厚重结构可以阻挡气流和声音传播，限制风元素机动。
    \item 暗克光：暗元素通过遮蔽和吸收削弱光元素的显现效果。
    \item 光克暗：光元素通过照明和辨识破除暗元素的隐匿效果。
\end{itemize}

注意，光与暗是双向限制关系。光可以照亮暗处，暗也可以吞没光线；哪一方占优，取决于强度、环境和使用者的术式结构。

\subsubsection{转化：改变表现形式}

\textbf{转化作用}指一种元素在特定条件下改变表现形式，转向另一种元素倾向。转化通常需要媒介、环境或第三种元素参与。

常见例子包括：
\begin{itemize}
    \item 水受寒可以转为偏土性的冰，表现出稳定和阻挡。
    \item 火与土结合可以形成熔岩，兼具流动、热量和实体破坏。
    \item 雷通过水传播时，攻击范围扩大，但控制难度上升。
    \item 光被暗元素封存后，可以形成延迟释放的封印术式。
\end{itemize}

为了描述元素效果，可以使用一个简单的设定公式：
\begin{equation}
    \text{最终效果}=\text{基础强度}\times\text{环境修正}\times\text{元素关系修正}
\end{equation}
该公式不用于精密计算，而是元素学中常见的简化表达：同一种术式在不同环境与不同元素关系下，最终效果可能明显不同。

\subsection{元素理论对世界架构的影响}\label{subsec:elementinfluence}

元素理论不仅用于解释个体能力，也用于解释地区、生态与文明差异。一个地区长期受到某种元素影响，它的地貌、产业、信仰和战斗方式都会发生变化。

\subsubsection{元素与地区}

不同元素富集的地区，会形成不同的自然和社会特征：
\begin{itemize}
    \item \textbf{土元素富集区}：多山脉、矿区、峡谷和城塞。居民擅长建筑、采矿、防御术和重型装备。
    \item \textbf{水元素富集区}：多河流、湖泊、湿地和港口。医疗、航运、药剂学和净化术较发达。
    \item \textbf{风元素富集区}：多高原、草原、峡风带和空港。通信、侦察、飞行器和游牧文化常见。
    \item \textbf{火元素富集区}：多火山、熔岩地、温泉和锻炉城市。锻造、能源、战争技术和祭火仪式发达。
    \item \textbf{雷元素富集区}：多风暴带、磁石山、机械都市和实验设施。机械文明、导能武器和神经强化技术常见。
    \item \textbf{光元素富集区}：多圣域、观测塔、学院和秩序城市。记录、审判、结界、教育和医疗制度完善。
    \item \textbf{暗元素富集区}：多地下城、边境、古墓、封印地和隐秘组织。潜行、保存、禁术、情报和梦境研究常见。
\end{itemize}

\subsubsection{元素亲和性}

\begin{definition}[元素亲和性]
    \textbf{元素亲和性（Elemental Affinity）}是个体、地区、种族或道具对某类元素的适应程度。亲和性越高，越容易感知、储存、引导或使用该元素。
\end{definition}

元素亲和性不是“只能使用某一种元素”。一个水亲和者仍然可以使用火元素，只是消耗更大、控制更难、稳定性更差。相反，高亲和性也不全是优点。火亲和者可能更容易受高温和情绪波动影响；暗亲和者在潜行方面更具优势，但在强光结界中更容易暴露异常。

亲和性可以来自多种来源：
\begin{itemize}
    \item \textbf{出生环境}：长期生活在某种元素富集地区的人，更容易形成对应亲和性。
    \item \textbf{血统与种族}：某些种族天生适应特定元素，如水生族群偏水，翼族偏风。
    \item \textbf{训练与契约}：术士可以通过修行、仪式或道具契约强化某种元素亲和。
    \item \textbf{创伤与污染}：强烈元素事故可能永久改变个体的亲和性，使其获得能力，也留下代价。
\end{itemize}

\subsubsection{元素与个体能力}

元素亲和性会影响个体更容易掌握哪类术式、在哪些环境中更稳定，以及受到哪类元素压制时更脆弱。术士学院通常会从四个方面记录个体的元素资料：
\begin{enumerate}
    \item \textbf{主亲和}：个体最容易感知和引导的元素。
    \item \textbf{副亲和}：会影响术式表现形式的次要元素。
    \item \textbf{代价}：使用元素时产生的负担，如体力消耗、情绪波动、环境依赖或身体损伤。
    \item \textbf{受制条件}：个体在何种元素环境或术式结构下更容易失去稳定性。
\end{enumerate}

例如，火与光亲和较高的人更容易掌握圣火术、净化术和审判术，但也可能因术式过亮、过热而伤及自身。水与暗亲和较高的人更容易掌握梦境治疗、记忆封存和沉眠术，但在强光结界或高温干燥环境中稳定性会下降。元素亲和性提供的是能力倾向，最终表现仍取决于训练、道具、环境和个人选择。

\subsubsection{元素与文明}

元素也会影响文明发展。火元素充足的地区容易发展锻造和能源技术；雷元素充足的地区容易发展机械和通信；光元素充足的地区容易发展记录、法律和教育；暗元素充足的地区则可能重视封印、情报、地下交通和长期保存。

因此，元素分布会自然改变地区的生活方式。例如，风元素城市常见高塔、吊桥、滑翔设备和远距信号塔；土元素王国常见坚固城墙、矿业贵族和重装军队；水元素群岛则常重视航海、医术、净化仪式和港口贸易。

\subsubsection{本节小结}

在\wkn{}世界内部，元素学提供了理解自然现象、术式能力、地区差异和文明形态的一套基础语言。土、水、风、火、雷、光、暗分别描述结构、流动、传递、释放、激活、显现与遮蔽等倾向。通过元素分布与元素亲和性，可以解释不同地区、族群与个体之间的差异。

\section{OC 变换理论}\label{sec:octheory}

本章系统阐述 OC 变换理论，涵盖连续与离散两类模型。
该理论阐述了现实世界个体特征向二次元原创角色（OC）特征的映射与转换机制，是世界观中设主与OC转换的理论基础。
\begin{introduction}
    \item 特征向量函数与特征向量序列
    \item 连续时间 OC 变换
    \item 离散时间 OC 变换
    \item 综合实例
    \item 特征向量函数到特征向量序列的采样理论
    \item 快速 OC 变换
\end{introduction}

\subsection{特征向量函数与特征向量序列}

\subsubsection{自设与OC}
在OC圈中，\textbf{设主（Creator）}是现实世界中的个体，\textbf{原创角色（Original Character，OC）}是二次元世界中的虚构角色。设主通过创造设定、绘制形象、编写故事等方式，赋予OC独特的特征。

然而，设主与OC之间并非一对一的简单映射关系：一个设主可能创造多个OC。但设主可以为自己创造\textbf{自设（Self-Insert OC）}，即在某种程度上代表设主自身的OC。自设通常承载着设主的理想化形象、情感寄托或自我表达。

因此，自设与设主之间存在更为紧密的特征关联，这种关联可以通过OC变换理论进行量化描述。换言之，\textbf{OC变换其实是设主特征向自设特征的映射过程}，所有OC的特征都与设主相关，但只有自设的特征才直接反映设主的属性。

对于一个设主而言，存在OC就能在数学上进行OC变换；但对于一个OC而言，只有当它是自设时，才在\wkn{}的世界观中物理实现OC变换。

\subsubsection{特征向量函数与特征向量序列的定义}

任何一个角色，无论是现实世界中的个体还是二次元世界中的虚构人物，都可以通过一组特征来完整描述。
这些特征构成了角色的“身份指纹”，能够唯一地表征该角色的核心属性。在现实世界中，这些特征可能包括生理
特征（如身高、体重）、心理特征（如性格倾向、情绪状态）、社会特征（如职业、人际关系）等；而在二次元世界
中，角色的特征则可能体现为更抽象的艺术属性（如萌属性、角色设定、剧情功能等）。

为了建立统一的数学框架来描述角色的特征演化，我们引入\textbf{特征向量（Feature Vector）}的概念。每个角色在任意时
刻的状态都可以用一个多维向量来表示，其中每个分量对应一个特定的特征属性。这种表示方法不仅能够捕捉角
色的静态特征，更能描述特征随时间变化的动态过程。

基于这一思想，我们定义以下统一的特征向量函数与特征向量序列框架。

\begin{definition}[特征向量函数]
    设现实世界的特征由 $N$ 个可量化属性描述，OC 的特征由 $M$ 个抽象属性描述。

    \textbf{现实特征向量函数（Real-world Feature Vector Function）}为 $N$ 维列向量：
    \begin{equation}
        \mathbf{r}(t)=\bigl[r_1(t),\,r_2(t),\,\dots,\,r_N(t)\bigr]^{\mathsf{T}},
    \end{equation}
    其中分量 $r_i(t)$ 为可量化的现实属性，随时间连续变化。

    \textbf{OC 特征向量函数（OC Feature Vector Function）}为 $M$ 维列向量：
    \begin{equation}
        \mathbf{q}(t)=\bigl[q_1(t),\,q_2(t),\,\dots,\,q_M(t)\bigr]^{\mathsf{T}},
    \end{equation}
    其中分量 $q_j(t)$ 表示 OC 的抽象属性，随时间连续变化。
\end{definition}
对连续信号以采样间隔 $T_s$ 取样，得到离散的序列，即为特征向量序列。
\begin{definition}[特征向量序列]
    \textbf{现实特征向量序列（Real-world Feature Vector Sequence）}为 $N$ 维列向量：
    \begin{equation}
        \mathbf{r}[n]=\bigl[r_1[n],\,r_2[n],\,\dots,\,r_N[n]\bigr]^{\mathsf{T}},\quad n\in\mathbb{Z},
    \end{equation}
    其中 $r_i[n]=r_i(nT_s)$ 为采样得到的离散序列分量。

    \textbf{OC 特征向量序列（OC Feature Vector Sequence）}为 $M$ 维列向量：
    \begin{equation}
        \mathbf{q}[n]=\bigl[q_1[n],\,q_2[n],\,\dots,\,q_M[n]\bigr]^{\mathsf{T}},\quad n\in\mathbb{Z},
    \end{equation}
    其中 $q_j[n]=q_j(nT_s)$ 为采样得到的离散序列分量。
\end{definition}

\begin{remark}
    这些函数或序列的分量 $r_i(t)$、$q_j(t)$、$r_i[n]$、$q_j[n]$ 均为实值函数或序列，且是\textbf{因果的}，即在 $t<0$ 或 $n<0$ 时为零。这一因果性设定反映了角色特征的演化具有\textbf{时间方向性}，特征的变化只能依赖于当前和过去的状态，而不能依赖于未来的信息。
\end{remark}


\subsection{连续时间 OC 变换}

在现实世界中，个体的特征随时间连续变化；二次元 OC 的特征亦然。
因此，连续时间 OC 变换是最自然的模型。

\subsubsection{时域描述}

\begin{theorem}[连续时间 OC 变换的时域描述]
    \textbf{连续时间 OC 变换（Continuous Time OC Transform）}由多输入多输出连续 LTI 系统实现，输入——输出关系用卷积积分给出
    \begin{equation}
        \mathbf{q}(t)=\int_{-\infty}^{\infty}\mathbf{H}(\tau)\mathbf{r}(t-\tau)\,\md\tau = \mathbf{H}(t)*\mathbf{r}(t),
        \label{eq:conv-cont}
    \end{equation}
    其中，
    \begin{equation}
        \mathbf{H}(\tau) = \begin{bmatrix}
            h_{11}(\tau) & h_{12}(\tau) & \cdots & h_{1N}(\tau) \\
            h_{21}(\tau) & h_{22}(\tau) & \cdots & h_{2N}(\tau) \\
            \vdots       & \vdots       & \ddots & \vdots       \\
            h_{M1}(\tau) & h_{M2}(\tau) & \cdots & h_{MN}(\tau)
        \end{bmatrix}
        \in\mathbb{R}^{M\times N}
    \end{equation}
    称为系统的 \textbf{连续OC冲激响应矩阵}（Continuous OC Impulse-Response Matrix），相应的系统称为\textbf{连续时间 OC 变换系统}（Continuous Time OC Transform System）。
\end{theorem}

\begin{remark}
    \begin{enumerate}
        \item \(\mathbf{H}(t)\) 是因果的。
              即 \(h_{ij}(t)=0\) 当 \(t<0\)，这反映了 OC 特征的生成只能依赖于现实特征的过去状态。
        \item 式 \eqref{eq:conv-cont} 表明，当前 OC 特征 $\mathbf{q}(t)$ 取决于现实特征 $\mathbf{r}(t)$ 的整个历史，$\mathbf{H}(\tau)$ 量化各历史时刻对当下的影响权重。
        \item 注意这里的卷积是沿时间变量 $t$ 进行的\textbf{一维时间卷积}，不是图像处理中的二维空间卷积。二维卷积通常在像素平面的两个空间坐标上滑动卷积核，用于图像滤波、边缘检测等；而式 \eqref{eq:conv-cont} 中的卷积描述的是历史时间上的累积影响。
        \item 对于输出向量 \(\mathbf{q}(t)\) 的第 \(i\) 个分量有
              \begin{equation}
                  q_i(t) = \int_{-\infty}^{\infty} \left[ \sum_{j=1}^{N} h_{ij}(\tau) r_j(t-\tau) \right] \md\tau
                  = \sum_{j=1}^{N} \int_{-\infty}^{\infty} h_{ij}(\tau) r_j(t-\tau) \, \md\tau
              \end{equation}
              即
              每个输出分量 \(q_i(t)\) 是所有输入分量 \(r_j(t)\) 与对应冲激响应函数 \(h_{ij}(\tau)\) 卷积的叠加：
              \begin{equation}
                  q_i(t) = \sum_{j=1}^{N} (h_{ij} * r_j)(t)
              \end{equation}
    \end{enumerate}
\end{remark}

\subsubsection{\texorpdfstring{$s$}{s} 域描述}

由于卷积运算在时域中计算复杂，我们引入拉普拉斯变换，将时域卷积转化为\(s\)域乘法，进而简化描述与分析。

\begin{definition}[拉普拉斯变换]
    对任意连续函数 $f(t)$，其拉普拉斯变换（也叫双边拉普拉斯变换）定义为
    \begin{equation}
        \tilde{F}(s)=\mathcal{L}\{f(t)\}=\int_{-\infty}^{\infty}\me^{-st}f(t)\,\md t,\quad s=\sigma+\mj\omega.
    \end{equation}
\end{definition}

\begin{theorem}[OC 变换的 $s$ 域描述]
    对 $\mathbf{r}(t)$、$\mathbf{q}(t)$ 及 $\mathbf{H}(\tau)$ 的每个元素分别施行拉普拉斯变换，记
    \begin{equation}
        \tilde{\mathbf{R}}(s)=\mathcal{L}\{\mathbf{r}(t)\},\quad
        \tilde{\mathbf{Q}}(s)=\mathcal{L}\{\mathbf{q}(t)\},\quad
        \tilde{\mathbf{H}}(s)=\mathcal{L}\{\mathbf{H}(\tau)\}.
    \end{equation}
    则在 $s$ 域有
    \begin{equation}
        \tilde{\mathbf{Q}}(s)=\tilde{\mathbf{H}}(s)\tilde{\mathbf{R}}(s),
    \end{equation}
    其中 $\tilde{\mathbf{H}}(s)$ 称为系统的 \textbf{连续OC传递函数矩阵}（Continuous OC Transfer-Function Matrix）。
\end{theorem}

\begin{remark}
    要能使用OC变换的$s$域描述，必须满足\(\mathbf{r}(t)\) 和 \(\mathbf{H}(t)\) 本身是因果的，且满足拉普拉斯变换的收敛条件。
    如果不是因果的，单边拉普拉斯变换会造成信号损失，引起错误的变换；如果不满足收敛条件，拉普拉斯变换将无法定义。
\end{remark}

为了方便地描述拉普拉斯变换，我们可以使用\textbf{零极点图}来表示传递函数矩阵 $\tilde{\mathbf{H}}(s)$ 的性质。
\begin{definition}[零极点图]
    对于传递函数矩阵 $\tilde{\mathbf{H}}(s)$ 的每个元素 $h_{ij}(s)$，其\textbf{零点（Zeros）}为使分子为零的复数 $s$ 值，\textbf{极点（Poles）}为使分母为零的复数 $s$ 值。
    \textbf{零极点图（Pole-Zero Plot）}是在复平面上绘制零点和极点的位置图。
\end{definition}
\begin{remark}
    对于矩阵值函数 $\mathbf{H}(t)$ 的拉普拉斯变换 $\tilde{\mathbf{H}}(s)$，其零极点图具有以下性质：

    \begin{enumerate}
        \item \textbf{极点来源}：$\tilde{\mathbf{H}}(s)$ 的极点来源于其各元素分母的公共零点。

        \item \textbf{收敛域的垂直边界}：在复平面上，收敛域由一条垂直边界线 $\operatorname{Re}(s) = \sigma_c$ 界定，其中 $\sigma_c$ 是 $\tilde{\mathbf{H}}(s)$ 最右侧极点的实部。

        \item \textbf{区域划分}：该垂直线将复平面分为两个区域：
              \begin{itemize}
                  \item 右侧（$\operatorname{Re}(s) > \sigma_c$）：$\tilde{\mathbf{H}}(s)$ 收敛且解析
                  \item 左侧（$\operatorname{Re}(s) < \sigma_c$）：$\tilde{\mathbf{H}}(s)$ 发散或包含极点
              \end{itemize}

        \item \textbf{矩阵整体性}：整个矩阵 $\tilde{\mathbf{H}}(s)$ 的收敛域由其所有元素中收敛域最窄的那个决定，即所有极点的最右侧实部。
    \end{enumerate}
\end{remark}

考虑一个2×2传递函数矩阵的例子：
\[
    \mathbf{H}(s) = \begin{bmatrix}
        \dfrac{s+1}{(s+2)(s+3)}   & \dfrac{1}{s+2}   \\[1em]
        \dfrac{s-1}{s^2 + 2s + 5} & \dfrac{s+2}{s+4}
    \end{bmatrix}
\]

该矩阵的零极点分布为：
\begin{itemize}
    \item $H_{11}(s)$: 零点 $s = -1$，极点 $s = -2, -3$
    \item $H_{12}(s)$: 无零点，极点 $s = -2$
    \item $H_{21}(s)$: 零点 $s = 1$，极点 $s = -1 \pm 2\mj$
    \item $H_{22}(s)$: 零点 $s = -2$，极点 $s = -4$
\end{itemize}

\begin{center}
    \begin{tikzpicture}[scale=1.2, >=Stealth]

        % 坐标轴
        \draw[->, thick] (-5, 0) -- (2, 0) node[right] {$\mathrm{Re}(s)$};
        \draw[->, thick] (0, -3) -- (0, 3) node[above] {$\mj\mathrm{Im}(s)$};

        % 网格（可选）
        % \draw[gray!20] (-5,-3) grid (2,3);

        % 实轴刻度
        \foreach \x in {-4,-3,-2,-1,1}
        \draw (\x,0.1) -- (\x,-0.1) node[below] {$\x$};

        % 虚轴刻度  
        \foreach \y in {-2,-1,1,2}
        \draw (0.1,\y) -- (-0.1,\y) node[left] {$\y \mj$};

        % 绘制极点 (poles) - 用 × 表示
        \draw[red, thick] (-2-0.12,-0.12) -- (-2+0.12,0.12);
        \draw[red, thick] (-2-0.12,0.12)  -- (-2+0.12,-0.12);
        % \node[below right, red] at (-2,0) {$-2$};

        \draw[red, thick] (-3-0.12,-0.12) -- (-3+0.12,0.12);
        \draw[red, thick] (-3-0.12,0.12)  -- (-3+0.12,-0.12);
        % \node[below left, red] at (-3,0) {$-3$};

        \draw[red, thick] (-4-0.12,-0.12) -- (-4+0.12,0.12);
        \draw[red, thick] (-4-0.12,0.12)  -- (-4+0.12,-0.12);
        % \node[below left, red] at (-4,0) {$-4$};

        \draw[red, thick] (-1-0.12,2-0.12) -- (-1+0.12,2+0.12);
        \draw[red, thick] (-1-0.12,2+0.12)  -- (-1+0.12,2-0.12);
        % \node[above left, red] at (-1,2) {$-1+2\mj$};

        \draw[red, thick] (-1-0.12,-2-0.12) -- (-1+0.12,-2+0.12);
        \draw[red, thick] (-1-0.12,-2+0.12)  -- (-1+0.12,-2-0.12);
        % \node[below left, red] at (-1,-2) {$-1-2\mj$};

        % 绘制零点 (zeros) - 用 ○ 表示（不显示数值标签）
        \draw[blue, thick] (-1,0) circle (2pt);
        \draw[blue, thick] (1,0) circle (2pt);
        \draw[blue, thick] (-2,0) circle (2pt);

        % 收敛边界（最右侧极点的实部为 -1）
        \draw[red, dashed, thick] (-1, -3) -- (-1, 3) node[pos=0.95, above left] {$\mathrm{Re}(s) = -1$};

        % 图例
        \begin{scope}[shift={(3,2)}, font=\scriptsize]
            \draw[red, thick] (-0.06,-0.06) -- (0.06,0.06);
            \draw[red, thick] (-0.06,0.06) -- (0.06,-0.06);
            \node[right] at (0.18,0) {极点 (\texttimes)};
            \draw[blue, thick] (0,-0.35) circle (2pt);
            \node[right] at (0.18,-0.35) {零点 ($\circ$)};
            \draw[red, dashed, thick] (-0.08,-0.7) -- (0.12,-0.7);
            \node[right] at (0.18,-0.7) {收敛边界};
        \end{scope}

    \end{tikzpicture}
\end{center}

收敛域：$\mathrm{Re}(s) > -1$（收敛边界右侧区域）

注意到在\(s = -2\)处出现了\textbf{零极点相消}现象，即传递函数矩阵的某些元素在该点既有零点又有极点，导致该点在整体传递函数中不表现为极点。这种现象在OC变换系统分析中具有重要意义，因为它可能影响系统的稳定性和响应特性。

\subsubsection{连续时间 OC 存在性定理}

每个人固有地拥有现实特征向量函数 $\mathbf{r}(t)$ 和自指的 OC 冲激响应矩阵 $\mathbf{H}(t)$。只有当这两个属性满足下述条件时，才能通过卷积实际生成合法的 OC 特征 $\mathbf{q}(t)$，这时才称“该人存在该OC”。

\begin{theorem}[连续时间 OC 存在性定理]
    在\wkn{}世界观的 OC 变换公理体系中，给定一个人的 $\mathbf{r}(t)$ 和 $\mathbf{H}(t)$，该人存在OC的\textbf{充分必要条件}是以下两组条件同时成立：
    \begin{enumerate}
        \item \textbf{现实特征向量函数的适应性}：$\mathbf{r}(t)$ 是因果的，允许拉普拉斯变换且与系统作用后不引发发散。
              \begin{enumerate}
                  \item $\mathbf{r}(t)=0$ 当 $t<0$。
                        等价的\(s\) 域表述为：$\tilde{\mathbf{R}}(s)$ 的收敛域包含某个右半平面 $\mathrm{Re}(s)>\sigma_r$。
                  \item \(\tilde{\mathbf{R}}(s)\)的收敛域包含虚轴。这是保证现实人设不崩坏的必要条件。
              \end{enumerate}
        \item \textbf{OC变换系统的正则性}：$\mathbf{H}(t)$ 是因果且稳定的。
              \begin{itemize}
                  \item \textbf{因果性}：$\mathbf{H}(t) = 0$ 对所有 $t < 0$ 成立
                  \item \textbf{稳定性}：$\mathbf{H}(t)$ 的每个元素 $h_{ij}(t)$ 绝对可积：
                        \begin{equation}
                            \int_0^\infty |h_{ij}(\tau)| \, \md\tau < \infty, \quad \forall i,j
                        \end{equation}
                  \item \(s\) 域必要表述：$\tilde{\mathbf{H}}(s)$ 的收敛域包含某个右半平面 $\mathrm{Re}(s)>\sigma_{\mathbf{H}}$且包含虚轴。
              \end{itemize}
    \end{enumerate}
\end{theorem}

\begin{remark}
    \begin{enumerate}
        \item 上述“充分必要条件”是\wkn{}世界观内部的 OC 存在性判定规则，而非现实信号处理理论中关于一般 LTI 系统输入输出合法性的完整分类。
        \item 当 $\mathbf{r}(t)$ 和 $\mathbf{H}(t)$ 满足上述条件时，卷积积分 \eqref{eq:conv-cont} 收敛，并生成合法的 OC 特征 $\mathbf{q}(t)$。
        \item 在\(s\)域，没有等价的因果性的描述，只能结合时域描述来判断。
              但稳定性可以通过 $\tilde{\mathbf{H}}(s)$ 的收敛域来描述。
        \item 现实特征向量函数 $\mathbf{r}(t)$ 的适应性保证了 OC 特征 $\mathbf{q}(t)$ 不会因为现实特征的异常而崩坏。
        \item OC 变换系统的正则性保证了 OC 特征 $\mathbf{q}(t)$ 的生成过程是稳定的，不会因为系统冲激响应矩阵 $\mathbf{H}(t)$ 的异常而导致 OC 特征的发散或不稳定。
        \item 当满足上述条件时，时域卷积积分和拉普拉斯变换都可以换成单边的，只考虑\(t>0\)的部分。即
              \begin{equation}
                  \mathbf{q}(t) = \int_0^\infty \mathbf{H}(\tau) \mathbf{r}(t-\tau) \, \md\tau
              \end{equation}
              拉普拉斯变换的定义式可以改成
              \begin{equation}
                  \tilde{F}(s)=\mathcal{L}\{f(t)\}=\int_0^\infty\me^{-st}f(t)\,\md t,\quad s=\sigma+\mj\omega.
              \end{equation}
        \item 从实际意义上理解：当一个设主“存在OC”时，无论是否是自设，必然满足上述条件。此时能够进行数学上的 OC 变换。
        \item 但在\wkn{}的世界观中，只有当该OC是自设时，才物理实现OC变换，即设主通过OC变换生成的OC特征$\mathbf{q}(t)$才能在二次元世界中实际存在。
    \end{enumerate}
\end{remark}

\subsubsection{连续时间 OC 反变换与可逆性}

\begin{definition}[连续时间 OC 可逆系统]
    对于一个由因果、稳定传递函数矩阵 $\tilde{\mathbf{H}}(s) \in \mathbb{C}^{M \times N}$ 描述的连续时间 OC 变换系统，如果存在一个同样因果、稳定的 LTI 系统 $\tilde{\mathbf{G}}(s) \in \mathbb{C}^{N \times M}$，使得在收敛域内满足：
    \begin{equation}
        \tilde{\mathbf{G}}(s) \tilde{\mathbf{H}}(s) = \mathbf{I}_N,
    \end{equation}
    则称原 OC 变换系统是 \textbf{可逆的}，并称 $\tilde{\mathbf{G}}(s)$ 为其 \textbf{逆系统}。
\end{definition}

\begin{definition}[连续时间 OC 反变换]
    对于一个可逆的连续时间 OC 变换系统，其 \textbf{反变换} 定义为由逆系统 $\tilde{\mathbf{G}}(s)$ 实现的变换过程。在时域中，该反变换由如下卷积积分给出：
    \begin{equation}
        \mathbf{r}(t) = \int_{0}^{\infty} \mathbf{G}(\tau) \mathbf{q}(t-\tau) \, \md\tau,
        \label{eq:inv-cont}
    \end{equation}
    其中 $\mathbf{G}(\tau) = \mathcal{L}^{-1}\{\tilde{\mathbf{G}}(s)\}$ 是逆系统的连续OC冲激响应矩阵。此过程可以从 OC 特征 $\mathbf{q}(t)$ 中唯一地恢复出现实描述 $\mathbf{r}(t)$。
\end{definition}

\begin{theorem}[连续时间 OC 反变换存在性判据]
    \label{the:inv-cont-existence}
    在\wkn{}世界观的 OC 变换公理体系中，一个由有理矩阵 $\tilde{\mathbf{H}}(s)\in\mathbb{C}^{M \times N}$ 描述的连续时间 OC 变换，其因果、稳定的左逆系统 $\tilde{\mathbf{G}}(s)\in\mathbb{C}^{N \times M}$ 存在，当且仅当以下条件同时成立：

    \begin{enumerate}
        \item \textbf{正向系统属性：} 原系统 $\tilde{\mathbf{H}}(s)$ 本身是因果且稳定的。

        \item \textbf{维度条件：} 现实特征的维度 $N$ 不超过 OC 特征的维度 $M$，即
              \begin{equation}
                  N \leq M.
              \end{equation}

        \item \textbf{列满秩条件：} 原系统 $\tilde{\mathbf{H}}(s)$ 在闭右半平面列满秩，即
              \begin{equation}
                  \operatorname{rank}(\tilde{\mathbf{H}}(s)) = N,\quad \forall s\in\mathbb{C}\ \text{且}\ \operatorname{Re}(s)\geq 0.
              \end{equation}

        \item \textbf{稳定左逆条件：} 存在一个有理矩阵 $\tilde{\mathbf{G}}(s)$ 满足 $\tilde{\mathbf{G}}(s)\tilde{\mathbf{H}}(s)=\mathbf{I}_N$，且 $\tilde{\mathbf{G}}(s)$ 的所有极点严格位于左半平面：
              \begin{equation}
                  \{\,s\in\mathbb{C}\mid s\ \text{为}\ \tilde{\mathbf{G}}(s)\ \text{的极点}\,\}
                  \subset
                  \{s\in\mathbb{C}:\mathrm{Re}(s)<0\}.
              \end{equation}
    \end{enumerate}
\end{theorem}

\begin{remark}
    \begin{enumerate}
        \item 只有当\(\tilde{\mathbf{H}}(s)\) 为因果且稳定的系统时，连续时间OC反变换才有意义。
        \item 当 $N > M$ 时，系统本质上是压缩的，从低维的 OC 特征无法唯一确定高维的现实特征，此时反变换不存在。
        \item 第二点和第三点排除了维度压缩以及闭右半平面内的秩亏损，是可稳定反演的必要限制。
        \item 第四点是世界观内部的可实现性公理：若形式逆系统含有右半平面极点，则反变换在代数上可写出，但不能作为稳定的 OC 反变换系统实现。
        \item 从实际意义上理解：对于设主，可以任意指定某样东西为自己的“OC”，即便这个“OC”不是人们通常意义上的角色。当该“OC”满足上述判据时，设主变换成自己的“OC”后能通过反变换唯一地恢复出自己的现实特征。
        \item 这里的“当且仅当”是\wkn{}世界观内部的判定公理；若将该模型视为现实控制理论中的一般 MIMO LTI 系统，还需要额外的系统论分析。
              这等价于\textbf{设主的OC足够独特}，不会导致反变换不存在，从而保证设主与其“OC”之间的特征映射是双向的。
        \item 同样地，只有当该“OC”是自设时，才能物理实现OC反变换。
    \end{enumerate}
\end{remark}

\begin{theorem}[形式逆系统的构造]
    当满足定理\ref{the:inv-cont-existence}时，逆系统可按如下方式构造，并按极点条件判断其是否属于稳定 OC 反变换系统：

    \begin{enumerate}
        \item \textbf{当 $M = N$ 时}：
              形式逆系统由矩阵求逆直接得到：
              \begin{equation}
                  \tilde{\mathbf{G}}(s) = \tilde{\mathbf{H}}(s)^{-1}.
              \end{equation}

        \item \textbf{当 $M > N$ 时}：
              一个常用的形式左逆由 Moore-Penrose 伪逆给出：
              \begin{equation}
                  \tilde{\mathbf{G}}(s)=\tilde{\mathbf{H}}^{\sharp}(s) = \left( \tilde{\mathbf{H}}(s)^\dagger \tilde{\mathbf{H}}(s) \right)^{-1} \tilde{\mathbf{H}}(s)^\dagger,
              \end{equation}
              其中 $^\dagger$ 表示共轭转置。
    \end{enumerate}
\end{theorem}

\begin{remark}
    上述伪逆是频域中的形式构造。在\wkn{}世界观中，若所得 $\tilde{\mathbf{G}}(s)$ 满足极点条件，则它被规定为可实现的 OC 反变换系统；若不满足，则只保留为形式反演而不具备稳定的物理实现。
\end{remark}

\subsection{离散时间 OC 变换}

在计算机系统的实现中，现实与 OC 特征均以离散序列形式出现，此时需借助离散时间 OC 变换。此外，对于OC的局部变换，离散模型更为合适。

\subsubsection{时域描述}

\begin{theorem}[离散时间 OC 变换的时域描述]
    \textbf{离散时间 OC 变换（Discrete Time OC Transform）}仍为多输入多输出 LTI 系统，其输入——输出关系由卷积和给出
    \begin{equation}
        \mathbf{q}[n]=\sum_{k=-\infty}^{\infty}\mathbf{H}[k]\mathbf{r}[n-k]=\mathbf{H}[n]*\mathbf{r}[n],
        \label{eq:conv-disc}
    \end{equation}
    其中，
    \begin{equation}
        \mathbf{H}[k] = \begin{bmatrix}
            h_{11}[k] & h_{12}[k] & \cdots & h_{1N}[k] \\
            h_{21}[k] & h_{22}[k] & \cdots & h_{2N}[k] \\
            \vdots    & \vdots    & \ddots & \vdots    \\
            h_{M1}[k] & h_{M2}[k] & \cdots & h_{MN}[k] \\
        \end{bmatrix}
    \end{equation}
    为系统的 \textbf{离散OC冲激响应矩阵}（Discrete OC Impulse-Response Matrix），相应的系统称为离散时间 OC 变换系统（Discrete Time OC Transform System）。
\end{theorem}

\begin{remark}
    式\eqref{eq:conv-disc}表明，当前 OC 特征 $\mathbf{q}[n]$ 取决于现实特征 $\mathbf{r}[n]$ 的整个历史，$\mathbf{H}[k]$ 量化各历史时刻对当下的影响权重。
    对于输出向量 \(\mathbf{q}[n]\) 的第 \(i\) 个分量有
    \begin{equation}
        q_i[n] = \sum_{k=-\infty}^{\infty} \left[ \sum_{j=1}^{N} h_{ij}[k] r_j[n-k] \right]
        = \sum_{j=1}^{N} \sum_{k=-\infty}^{\infty} h_{ij}[k] r_j[n-k]
    \end{equation}
    即
    每个输出分量 \(q_i[n]\) 是所有输入分量 \(r_j[n]\) 与对应离散冲激响应函数 \(h_{ij}[k]\) 卷积的叠加：
    \begin{equation}
        q_i[n] = \sum_{j=1}^{N} (h_{ij} * r_j)[n]
    \end{equation}
\end{remark}

\subsubsection{\texorpdfstring{$z$}{z} 域描述}
和连续时间类似，离散时间的卷积和在时域中计算复杂，我们引入 $Z$ 变换，将时域卷积转化为\(z\)域乘法，进而简化描述与分析。

\begin{definition}[$Z$ 变换]
    对任意离散序列 $f[n]$，其 \(Z\) 变换定义为
    \begin{equation}
        \tilde{F}(z)=\mathcal{Z}\{f[n]\}=\sum_{n=-\infty}^{\infty}f[n]z^{-n},\quad z=A\me^{\mj\omega}.
    \end{equation}
\end{definition}

\begin{theorem}[OC 变换的 $z$ 域描述]
    对 $\mathbf{r}[n]$、$\mathbf{q}[n]$ 及 $\mathbf{H}[k]$ 的每个元素分别施行 $Z$ 变换，记
    \begin{equation}
        \tilde{\mathbf{R}}(z)=\mathcal{Z}\{\mathbf{r}[n]\},\quad
        \tilde{\mathbf{Q}}(z)=\mathcal{Z}\{\mathbf{q}[n]\},\quad
        \tilde{\mathbf{H}}(z)=\mathcal{Z}\{\mathbf{H}[k]\}.
    \end{equation}
    则在 $z$ 域有
    \begin{equation}
        \tilde{\mathbf{Q}}(z)=\tilde{\mathbf{H}}(z)\tilde{\mathbf{R}}(z),
    \end{equation}
    其中 $\tilde{\mathbf{H}}(z)$ 称为系统的 \textbf{离散OC传递函数矩阵}（Discrete OC Transfer-Function Matrix）。
\end{theorem}

与连续时间情况类似，我们可以使用\textbf{零极点图}来表示传递函数矩阵 $\tilde{\mathbf{H}}(z)$ 的性质。
\begin{definition}[零极点图（$z$ 域）]
    对于离散传递函数矩阵 $\tilde{\mathbf{H}}(z)$ 的每个元素 $h_{ij}(z)$，其\textbf{零点（Zeros）}为使分子为零的复数 $z$ 值，\textbf{极点（Poles）}为使分母为零的复数 $z$ 值。
    $z$ 平面上的\textbf{零极点图（Pole-Zero Plot）}是绘制这些零点和极点位置的图示。
\end{definition}

\begin{remark}
    对于矩阵值序列 $\mathbf{H}[n]$ 的 $Z$ 变换 $\tilde{\mathbf{H}}(z)$，其零极点图具有以下关键性质：

    \begin{enumerate}
        \item \textbf{极点来源}：$\tilde{\mathbf{H}}(z)$ 的极点来源于其各元素分母的公共零点。

        \item \textbf{收敛域的环形边界}：在 $z$ 平面上，收敛域由圆形边界 $|z| = R_c$ 界定，其中 $R_c$ 是 $\tilde{\mathbf{H}}(z)$ 最外侧极点的模值。

        \item \textbf{区域划分}：该圆将 $z$ 平面分为两个区域：
              \begin{itemize}
                  \item 圆外区域（$|z| > R_c$）：$\tilde{\mathbf{H}}(z)$ 收敛且解析
                  \item 圆内区域（$|z| < R_c$）：$\tilde{\mathbf{H}}(z)$ 发散或包含极点
              \end{itemize}

              % \item \textbf{稳定性判据}：对于因果系统，\textbf{所有极点必须位于单位圆内（$|z| < 1$）}才能保证系统稳定。这是离散时间系统与连续时间系统最重要的区别之一。

        \item \textbf{矩阵整体性}：整个矩阵 $\tilde{\mathbf{H}}(z)$ 的收敛域由其所有元素中收敛域最窄的那个决定，即所有极点的最大模值。
    \end{enumerate}
\end{remark}

考虑一个2×2离散传递函数矩阵的例子：
\[
    \mathbf{H}(z) = \begin{bmatrix}
        \dfrac{z+0.5}{(z-0.2)(z+0.8)}   & \dfrac{1}{z-0.8}     \\[1em]
        \dfrac{z-0.3}{z^2 - 1.2z + 0.5} & \dfrac{z+0.8}{z-0.9}
    \end{bmatrix}
\]

该矩阵的零极点分布为：
\begin{itemize}
    \item $H_{11}(z)$: 零点 $z = -0.5$，极点 $z = 0.2, -0.8$
    \item $H_{12}(z)$: 无零点，极点 $z = 0.8$
    \item $H_{21}(z)$: 零点 $z = 0.3$，极点 $z = 0.6 \pm 0.447\mj$
    \item $H_{22}(z)$: 零点 $z = -0.8$，极点 $z = 0.9$
\end{itemize}

\begin{center}
    \begin{tikzpicture}[
            scale=1.5,                     % 1. 图大一点
            >=Stealth,
            font=\footnotesize,
            pole/.style={red, line width=.8pt},
            zero/.style={blue, line width=.8pt}
        ]

        % ---------------- 坐标轴 ----------------
        \draw[->, thick] (-2.4,0) -- (2.4,0) node[right] {$\mathrm{Re}(z)$};
        \draw[->, thick] (0,-2.4) -- (0,2.4) node[above] {$\mathrm{Im}(z)$};

        % ---------------- 单位圆 ----------------
        \def\Runit{1}
        \draw[blue, very thick] (0,0) circle (\Runit);
        % \node[blue, above=1pt] at (90:\Runit) {单位圆};

        % ---------------- 刻度 ----------------
        \foreach \x in {-1,1}{
                \draw (\x,0.05) -- (\x,-0.05) node[below=1pt] {$\x$};
                \draw (0,\x)  circle (0.05) node[left=1pt]  {$\x\mathrm{j}$};
            }

        % ---------------- 极点（仅画 ×，不标值） ----------------
        \def\cross#1#2{
            \draw[pole] (#1-0.06,#2-0.06) -- (#1+0.06,#2+0.06)
            (#1-0.06,#2+0.06) -- (#1+0.06,#2-0.06);
        }
        \cross{0.2}{0}
        \cross{-0.8}{0}
        \cross{0.9}{0}
        \cross{0.6}{0.447}
        \cross{0.6}{-0.447}
        \cross{0.8}{0}
        % \cross{0}{0}

        % ---------------- 零点（仅画 ○，不标值） ----------------
        \def\zerocircle#1#2{\draw[zero] (#1,#2) circle (2pt);}
        \zerocircle{-0.5}{0}
        \zerocircle{0.3}{0}
        \zerocircle{-0.8}{0}
        % \zerocircle{0}{0}  % H12无零点，此处画一个空心圆表示无零点

        % ---------------- 收敛边界（标注改进） ----------------
        \def\Rroc{0.9}
        \draw[red, dashed, thick] (0,0) circle (\Rroc);
        % 标签放在圆外侧并带白底以避免与图形重叠
        \node[red, font=\small, fill=white, inner sep=1pt] at (45:{\Rroc+0.08}) {$|z|=0.9$};

        % ---------------- 图例 ----------------
        \begin{scope}[shift={(2.1,2.0)}, font=\scriptsize]
            \draw[pole] (-0.06,-0.06) -- (0.06,0.06);
            \draw[pole] (-0.06,0.06) -- (0.06,-0.06);
            \node[right] at (0.18,0) {极点 (\texttimes)};
            \draw[zero] (0,-0.35) circle (2pt);
            \node[right] at (0.18,-0.35) {零点 ($\circ$)};
            \draw[red, dashed, thick] (-0.08,-0.7) -- (0.12,-0.7);
            \node[right] at (0.18,-0.7) {收敛边界};
            \draw[blue, very thick] (-0.08,-1.05) -- (0.12,-1.05);
            \node[right] at (0.18,-1.05) {单位圆};
        \end{scope}

    \end{tikzpicture}
\end{center}

收敛域：$|z| > 0.9$（收敛圆外区域），包含\(z = \infty\)。

注意到在\(z = 0.8\)处出现了\textbf{零极点相消}现象。在离散时间 OC 变换系统中，这种现象同样会影响系统的稳定性和响应特性。特别需要注意的是，由于存在模值为 $0.9$ 的极点（位于单位圆内但接近边界），该系统的稳定性处于临界状态，其 OC 特征的生成过程可能会表现出较慢的衰减或持续的振荡。

\subsubsection{离散时间 OC 存在性定理}

每个人固有地拥有现实特征向量序列 $\mathbf{r}[n]$ 和自指的 OC 冲激响应矩阵 $\mathbf{H}[n]$。只有当这两个属性满足下述条件时，才能通过卷积和实际生成合法的 OC 特征 $\mathbf{q}[n]$，这时才称“该人存在该OC”。

\begin{theorem}[离散时间 OC 存在性定理]
    在\wkn{}世界观的 OC 变换公理体系中，给定一个人的 $\mathbf{r}[n]$ 和 $\mathbf{H}[n]$，该人存在OC的\textbf{充分必要条件}是以下两组条件同时成立：
    \begin{enumerate}
        \item \textbf{现实特征向量序列的适应性}：$\mathbf{r}[n]$ 是因果的，允许 $Z$ 变换且与系统作用后不引发发散。
              \begin{enumerate}
                  \item $\mathbf{r}[n]=0$ 当 $n<0$。
                        等价的$z$ 域表述为：$\tilde{\mathbf{R}}(z)$ 的收敛域包含某个圆外区域 $|z|>R_r$ 并包含 \(z=\infty\)。
                  \item \(\tilde{\mathbf{R}}(z)\)的收敛域包含单位圆。这是保证现实人设不崩坏的必要条件。
              \end{enumerate}
        \item \textbf{OC变换系统的正则性}：$\mathbf{H}[n]$ 是因果且稳定的。
              \begin{itemize}
                  \item \textbf{因果性}：$\mathbf{H}[n] = 0$ 对所有 $n < 0$ 成立
                  \item \textbf{稳定性}：$\mathbf{H}[n]$ 的每个元素 $h_{ij}[n]$ 绝对可和：
                        \begin{equation}
                            \sum_{n=0}^\infty |h_{ij}[n]| < \infty, \quad \forall i,j
                        \end{equation}
                  \item \textbf{$z$域等价表述}：$\tilde{\mathbf{H}}(z)$ 的收敛域为某个圆外区域 $|z| > R_{\text{max}}$，且包含单位圆 $|z| = 1$ 和无穷大 \(z=\infty\)。
              \end{itemize}
    \end{enumerate}
\end{theorem}

\begin{remark}
    \begin{enumerate}
        \item 上述“充分必要条件”是\wkn{}世界观内部的 OC 存在性判定规则，而非现实离散信号处理理论中关于一般 LTI 系统输入输出合法性的完整分类。
        \item 与连续时间OC变换不同，\(z\)域描述存在因果性的等价表述，即收敛域包含\(z=\infty\)。
        \item 当满足上述条件时，时域卷积和和\(Z\)变换都可以换成单边的，即
              \begin{equation}
                  \mathbf{q}[n] = \sum_{k=0}^{\infty} \mathbf{H}[k] \mathbf{r}[n-k]
              \end{equation}
              \(Z\)变换的定义式可以改成
              \begin{equation}
                  \tilde{F}(z)=\mathcal{Z}\{f[n]\}=\sum_{n=0}^{\infty}f[n]z^{-n},\quad z=A\me^{\mj\omega}.
              \end{equation}
    \end{enumerate}
\end{remark}

\subsubsection{离散时间 OC 反变换与可逆性}

\begin{definition}[离散时间 OC 可逆系统]
    对于一个由因果、稳定传递函数矩阵 $\tilde{\mathbf{H}}(z) \in \mathbb{C}^{M \times N}$ 描述的离散时间 OC 变换系统，如果存在一个同样因果、稳定的 LTI 系统 $\tilde{\mathbf{G}}(z) \in \mathbb{C}^{N \times M}$，使得在收敛域内满足：
    \begin{equation}
        \tilde{\mathbf{G}}(z) \tilde{\mathbf{H}}(z) = \mathbf{I}_N,
    \end{equation}
    则称原 OC 变换系统是 \textbf{可逆的}，并称 $\tilde{\mathbf{G}}(z)$ 为其 \textbf{逆系统}。
\end{definition}

\begin{definition}[离散时间 OC 反变换]
    对于一个可逆的离散时间 OC 变换系统，其 \textbf{反变换} 定义为由逆系统 $\tilde{\mathbf{G}}(z)$ 实现的变换过程。在时域中，该反变换由如下卷积和给出：
    \begin{equation}
        \mathbf{r}[n] = \sum_{k=0}^{\infty} \mathbf{G}[k] \mathbf{q}[n-k],
        \label{eq:inv-disc}
    \end{equation}
    其中 $\mathbf{G}[k] = \mathcal{Z}^{-1}\{\tilde{\mathbf{G}}(z)\}$ 是逆系统的离散OC冲激响应矩阵。此过程可以从 OC 特征 $\mathbf{q}[n]$ 中唯一地恢复出现实描述 $\mathbf{r}[n]$。
\end{definition}

\begin{theorem}[离散时间 OC 反变换存在性判据]
    \label{the:inv-disc-existence}
    在\wkn{}世界观的 OC 变换公理体系中，一个由有理矩阵 $\tilde{\mathbf{H}}(z)\in\mathbb{C}^{M \times N}$ 描述的离散时间 OC 变换，其因果、稳定的左逆系统 $\tilde{\mathbf{G}}(z)\in\mathbb{C}^{N \times M}$ 存在，当且仅当以下条件同时成立：

    \begin{enumerate}
        \item \textbf{正向系统属性：} 原系统 $\tilde{\mathbf{H}}(z)$ 本身是因果且稳定的。

        \item \textbf{维度条件：} 现实特征的维度 $N$ 不超过 OC 特征的维度 $M$，即
              \begin{equation}
                  N \leq M.
              \end{equation}

        \item \textbf{满秩条件：} 原系统 $\tilde{\mathbf{H}}(z)$ 列满秩，即
              \begin{equation}
                  \operatorname{rank}(\tilde{\mathbf{H}}(z)) = N,\quad \forall z\in\mathbb{C}\ \text{且}\ |z|\geq 1.
              \end{equation}
              其等价的表述为：$\tilde{\mathbf{H}}(z)$ 的各个列向量线性无关，或$\tilde{\mathbf{H}}(z)$ 作为线性映射的矩阵时，对应的映射是单射。

        \item \textbf{稳定左逆条件：} 存在一个有理矩阵 $\tilde{\mathbf{G}}(z)$ 满足 $\tilde{\mathbf{G}}(z)\tilde{\mathbf{H}}(z)=\mathbf{I}_N$，且 $\tilde{\mathbf{G}}(z)$ 的所有极点严格位于单位圆内：
              \begin{equation}
                  \{\,z\in\mathbb{C}\mid z\ \text{为}\ \tilde{\mathbf{G}}(z)\ \text{的极点}\,\}
                  \subset
                  \{z\in\mathbb{C}:|z|<1\}.
              \end{equation}
    \end{enumerate}
\end{theorem}

\begin{remark}
    这里的“当且仅当”是\wkn{}世界观内部的判定公理。若将该模型视为现实离散控制理论中的一般 MIMO LTI 系统，还需要额外分析稳定左逆、零点与可实现性等问题。
\end{remark}

\begin{theorem}[形式逆系统的构造]
    当满足定理\ref{the:inv-disc-existence}时，逆系统可按如下方式构造，并按极点条件判断其是否属于稳定 OC 反变换系统：

    \begin{enumerate}
        \item \textbf{当 $M = N$ 时}：
              形式逆系统由矩阵求逆直接得到：
              \begin{equation}
                  \tilde{\mathbf{G}}(z) = \tilde{\mathbf{H}}(z)^{-1}.
              \end{equation}

        \item \textbf{当 $M > N$ 时}：
              一个常用的形式左逆由 Moore-Penrose 伪逆给出：
              \begin{equation}
                  \tilde{\mathbf{G}}(z)=\tilde{\mathbf{H}}^{\sharp}(z) = \left( \tilde{\mathbf{H}}(z)^\dagger \tilde{\mathbf{H}}(z) \right)^{-1} \tilde{\mathbf{H}}(z)^\dagger,
              \end{equation}
              其中 $^\dagger$ 表示共轭转置。
    \end{enumerate}
\end{theorem}

\subsection{综合实例}
\begin{example}[连续时间OC变换综合题]

    某画师擅长为他人设计OC，每个人的现实特征向量函数的特征维度 $N=2$，画师设计的OC的特征维度 $M=3$。某个人的OC的变换系统的冲激响应矩阵为：
    \begin{equation}
        \mathbf{H}(t) = \begin{bmatrix}
            2\me^{-3t} - \me^{-2t} & \me^{-3t}              \\
            \me^{-2t}              & 2\me^{-3t} - \me^{-2t} \\
            \me^{-3t}              & \me^{-2t}
        \end{bmatrix} \cdot u(t)
    \end{equation}
    其中 $u(t)$ 为单位阶跃函数。

    该人的现实特征向量函数为：
    \begin{equation}
        \mathbf{r}(t) = \begin{bmatrix}
            \me^{-t} \\ 2\me^{-2t}
        \end{bmatrix} \cdot u(t)
    \end{equation}

    请回答以下问题：
    \begin{enumerate}
        \item 验证该OC变换系统是否满足连续时间OC存在性定理的所有条件。
        \item 求系统的连续OC传递函数矩阵 $\tilde{\mathbf{H}}(s)$，并确定其收敛域。
        \item 计算OC特征向量函数 $\mathbf{q}(t)$ 的时域表达式。
        \item 判断该系统是否存在因果稳定的逆系统，并说明理由。
        \item 如果存在逆系统，求其传递函数矩阵 $\tilde{\mathbf{G}}(s)$。
    \end{enumerate}

    \solution
    \begin{enumerate}
        \item \textbf{存在性验证：}
              \begin{itemize}
                  \item \textbf{现实特征向量函数的适应性：} $\mathbf{r}(t)$ 是因果的（$t<0$ 时为0），且每个分量拉普拉斯变换存在：
                        $\mathcal{L}\{\me^{-t}\} = \dfrac{1}{s+1}$（收敛域 $\mathrm{Re}(s)>-1$），
                        $\mathcal{L}\{2\me^{-2t}\} = \dfrac{2}{s+2}$（收敛域 $\mathrm{Re}(s)>-2$）。
                  \item \textbf{OC变换系统的正则性：} $\mathbf{H}(t)$ 是因果的（含 $u(t)$），且每个元素绝对可积：
                        如 $\displaystyle\int_0^\infty |2\me^{-3t}-\me^{-2t}|\,\md t < \infty$ 等，系统稳定。
              \end{itemize}
              满足所有存在性条件。

        \item \textbf{传递函数矩阵：}
              利用 $\mathcal{L}\{\me^{-at}u(t)\} = \dfrac{1}{s+a}$，可得：
              \begin{equation}
                  \tilde{\mathbf{H}}(s) = \begin{bmatrix}
                      \dfrac{2}{s+3} - \dfrac{1}{s+2} & \dfrac{1}{s+3}                  \\
                      \dfrac{1}{s+2}                  & \dfrac{2}{s+3} - \dfrac{1}{s+2} \\
                      \dfrac{1}{s+3}                  & \dfrac{1}{s+2}
                  \end{bmatrix}
                  = \begin{bmatrix}
                      \dfrac{s+1}{(s+2)(s+3)} & \dfrac{1}{s+3}          \\
                      \dfrac{1}{s+2}          & \dfrac{s+1}{(s+2)(s+3)} \\
                      \dfrac{1}{s+3}          & \dfrac{1}{s+2}
                  \end{bmatrix}
              \end{equation}
              收敛域为 $\mathrm{Re}(s) > -2$（包含虚轴，故稳定）。

        \item \textbf{OC特征向量函数：}
              在 $s$ 域计算 $\tilde{\mathbf{Q}}(s) = \tilde{\mathbf{H}}(s)\tilde{\mathbf{R}}(s)$：
              \begin{equation}
                  \tilde{\mathbf{R}}(s) = \begin{bmatrix}
                      \dfrac{1}{s+1} \\ \dfrac{2}{s+2}
                  \end{bmatrix}
              \end{equation}
              \begin{equation}
                  \tilde{\mathbf{Q}}(s) = \begin{bmatrix}
                      \dfrac{s+1}{(s+2)(s+3)}\cdot\dfrac{1}{s+1} + \dfrac{1}{s+3}\cdot\dfrac{2}{s+2} \\
                      \dfrac{1}{s+2}\cdot\dfrac{1}{s+1} + \dfrac{s+1}{(s+2)(s+3)}\cdot\dfrac{2}{s+2} \\
                      \dfrac{1}{s+3}\cdot\dfrac{1}{s+1} + \dfrac{1}{s+2}\cdot\dfrac{2}{s+2}
                  \end{bmatrix}
                  = \begin{bmatrix}
                      \dfrac{3}{(s+2)(s+3)}                                \\
                      \dfrac{1}{(s+1)(s+2)} + \dfrac{2(s+1)}{(s+2)^2(s+3)} \\
                      \dfrac{1}{(s+1)(s+3)} + \dfrac{2}{(s+2)^2}
                  \end{bmatrix}
              \end{equation}
              反拉普拉斯变换得：
              \begin{equation}
                  \mathbf{q}(t) = \begin{bmatrix}
                      3(\me^{-2t} - \me^{-3t})                                                                \\
                      (\me^{-t} - \me^{-2t}) + 2\!\left(\dfrac{1}{2}t\me^{-2t} + \me^{-2t} - \me^{-3t}\right) \\
                      \dfrac{1}{2}(\me^{-t} - \me^{-3t}) + 2t\me^{-2t}
                  \end{bmatrix} u(t)
              \end{equation}

        \item \textbf{逆系统存在性：} 不存在因果稳定的逆系统。理由：
              \begin{itemize}
                  \item 维度条件 $N=2 \leq M=3$ 满足。
                  \item 为判断列满秩性，只需考察两个列向量是否可能线性相关。取第一、二行构成的 $2\times2$ 子式：
                        \begin{equation}
                            \Delta(s)=
                            \begin{vmatrix}
                                \dfrac{s+1}{(s+2)(s+3)} & \dfrac{1}{s+3} \\
                                \dfrac{1}{s+2}          & \dfrac{s+1}{(s+2)(s+3)}
                            \end{vmatrix}
                            =\dfrac{(s+1)^2-(s+2)(s+3)}{(s+2)^2(s+3)^2}
                            =\dfrac{-3s-5}{(s+2)^2(s+3)^2}。
                        \end{equation}
                        该子式在 $s=-\dfrac{5}{3}$ 处为零，但此点位于左半平面，不影响闭右半平面的列满秩要求。要完全判定逆系统是否稳定，还需检查形式左逆的极点位置。
                  \item 因此，仅凭本题给出的简要计算不能断言逆系统不存在；在保守判据下，可以说该系统满足维度条件与闭右半平面的列满秩检查，但仍需进一步构造并检查形式左逆的稳定性。
              \end{itemize}

        \item 由于本题的核心是说明判据流程，完整形式左逆的构造可按定理~\ref{the:inv-cont-existence} 后的形式逆系统方法进行，并检查所得极点是否位于左半平面。
    \end{enumerate}
\end{example}

\begin{example}[离散时间OC变换综合题]

    某画师和计算机专业的助手共同创建了一个离散时间OC变换系统，能接受现实特征维度 $N=2$的现实特征向量序列，输出的OC特征向量序列的特征维度 $M=3$。系统的离散OC冲激响应矩阵为：
    \begin{equation}
        \mathbf{H}[n] = \begin{bmatrix}
            \left(\dfrac{1}{2}\right)^n - \left(\dfrac{1}{3}\right)^n & \left(\dfrac{1}{3}\right)^n                               \\
            \left(\dfrac{1}{2}\right)^n                               & \left(\dfrac{1}{2}\right)^n - \left(\dfrac{1}{3}\right)^n \\
            \left(\dfrac{1}{3}\right)^n                               & \left(\dfrac{1}{2}\right)^n
        \end{bmatrix} \cdot u[n]
    \end{equation}
    其中 $u[n]$ 为单位阶跃序列。

    该创作者在 $n \geq 0$ 时的现实特征向量序列为：
    \begin{equation}
        \mathbf{r}[n] = \begin{bmatrix}
            \left(\dfrac{1}{4}\right)^n \\ 2\left(\dfrac{1}{2}\right)^n
        \end{bmatrix} \cdot u[n]
    \end{equation}

    请回答以下问题：
    \begin{enumerate}
        \item 验证该OC变换系统是否满足离散时间OC存在性定理的所有条件。
        \item 求系统的离散OC传递函数矩阵 $\tilde{\mathbf{H}}(z)$，并确定其收敛域。
        \item 计算OC特征向量序列 $\mathbf{q}[n]$ 的时域表达式。
        \item 判断该系统是否存在因果稳定的逆系统，并说明理由。
        \item 如果存在逆系统，证明 $\tilde{\mathbf{H}}^\dagger(z)\tilde{\mathbf{H}}(z)$ 是可逆的，并给出其逆矩阵的表达式。
    \end{enumerate}

    \solution
    \begin{enumerate}
        \item \textbf{存在性验证：}
              \begin{itemize}
                  \item \textbf{现实特征向量序列的适应性：} $\mathbf{r}[n]$ 是因果的，且每个分量Z变换存在：
                        $\mathcal{Z}\{(1/4)^n\} = \dfrac{z}{z-1/4}$（收敛域 $|z|>1/4$），
                        $\mathcal{Z}\{2(1/2)^n\} = \dfrac{2z}{z-1/2}$（收敛域 $|z|>1/2$）。
                  \item \textbf{OC变换系统的正则性：} $\mathbf{H}[n]$ 是因果的，且每个元素绝对可和：
                        如 $\displaystyle\sum_{n=0}^\infty \left|\left(\dfrac{1}{2}\right)^n - \left(\dfrac{1}{3}\right)^n\right| < \infty$ 等，系统稳定。
              \end{itemize}
              满足所有存在性条件。

        \item \textbf{传递函数矩阵：}
              利用 $\mathcal{Z}\{a^n u[n]\} = \dfrac{z}{z-a}$，可得：
              \begin{equation}
                  \tilde{\mathbf{H}}(z) = \begin{bmatrix}
                      \dfrac{z}{z-1/2} - \dfrac{z}{z-1/3} & \dfrac{z}{z-1/3}                    \\
                      \dfrac{z}{z-1/2}                    & \dfrac{z}{z-1/2} - \dfrac{z}{z-1/3} \\
                      \dfrac{z}{z-1/3}                    & \dfrac{z}{z-1/2}
                  \end{bmatrix}
                  = \begin{bmatrix}
                      \dfrac{z/6}{(z-1/2)(z-1/3)} & \dfrac{z}{z-1/3}            \\
                      \dfrac{z}{z-1/2}            & \dfrac{z/6}{(z-1/2)(z-1/3)} \\
                      \dfrac{z}{z-1/3}            & \dfrac{z}{z-1/2}
                  \end{bmatrix}
              \end{equation}
              收敛域为 $|z| > 1/2$（包含单位圆，故稳定）。

        \item \textbf{OC特征向量序列：}
              在 $z$ 域计算 $\tilde{\mathbf{Q}}(z) = \tilde{\mathbf{H}}(z)\tilde{\mathbf{R}}(z)$：
              \begin{equation}
                  \tilde{\mathbf{R}}(z) = \begin{bmatrix}
                      \dfrac{z}{z-1/4} \\ \dfrac{2z}{z-1/2}
                  \end{bmatrix}
              \end{equation}
              以第一分量为例：
              \begin{equation}
                  \tilde{Q}_1(z) = \dfrac{z/6}{(z-1/2)(z-1/3)}\cdot\dfrac{z}{z-1/4} + \dfrac{z}{z-1/3}\cdot\dfrac{2z}{z-1/2}
              \end{equation}
              反Z变换得各分量表达式（具体计算略）。

        \item \textbf{逆系统存在性：} 存在因果稳定的逆系统。理由：
              \begin{itemize}
                  \item 维度条件 $N=2 \leq M=3$ 满足
                  \item $\tilde{\mathbf{H}}(z)$ 在所关注的收敛域 $|z|\geq 1$ 对应位置列满秩
                  \item 因此可构造因果稳定的逆系统（例如通过伪逆构造）
              \end{itemize}

        \item \textbf{伪逆验证:}
              \begin{equation}
                  \tilde{\mathbf{H}}^\dagger(z)\tilde{\mathbf{H}}(z) = \begin{bmatrix}
                      A(z) & B(z) \\ B(z) & A(z)
                  \end{bmatrix}
              \end{equation}
              其中
              \begin{equation}
                  A(z) = \left|\dfrac{z}{z-1/2}\right|^2 + \left|\dfrac{z}{z-1/3}\right|^2 + \left|\dfrac{z}{z-1/3}\right|^2,
              \end{equation}
              \begin{equation}
                  B(z) = \dfrac{z/6}{(z-1/2)(z-1/3)}\cdot\dfrac{\overline{z}}{\overline{z}-1/3} + \dfrac{\overline{z}}{\overline{z}-1/2}\cdot\dfrac{z}{z-1/3} + \dfrac{\overline{z}}{\overline{z}-1/3}\cdot\dfrac{z}{z-1/2}.
              \end{equation}

              该矩阵在 $|z|\geq 1$ 时满秩，故可逆。其逆矩阵为：
              \begin{equation}
                  [\tilde{\mathbf{H}}^\dagger(z)\tilde{\mathbf{H}}(z)]^{-1} = \dfrac{1}{A(z)^2 - B(z)^2} \begin{bmatrix}
                      A(z) & -B(z) \\ -B(z) & A(z)
                  \end{bmatrix}.
              \end{equation}
    \end{enumerate}
\end{example}

\subsection{特征向量函数到特征向量序列的采样理论}
\label{subsec:sampling-theory}

当现实世界的设主要穿越到二次元世界成为其自设，或自设从二次元世界返回到现实世界时，均需要进行OC变换或反变换。
这一过程包括：
\begin{enumerate}
    \item 从现实世界的特征向量函数到OC特征向量函数的变化，此时发生了连续时间OC变换；
    \item 从OC特征向量函数到现实世界的特征向量函数的变化，此时发生了连续时间OC反变换。
\end{enumerate}

在计算连续时间OC变换与反变换的步骤中，如果人为手算显然不可能。基于此，OC设计者、画师、计算机专家等人会共同设计基于现代计算机的\textbf{自动OC变换机}。
自动OC变换机无法直接处理连续的特征向量函数，只能处理离散的特征向量序列。
因此，自动OC变换机的实现需要结合采样理论，将连续时间OC变换离散化为离散时间OC变换。


\subsubsection{采样过程与频谱分析}

\begin{definition}[OC变换的采样过程]
    对连续时间现实特征向量函数$\mathbf{r}(t)$以采样间隔$T_s$进行均匀采样，得到离散序列：
    \begin{equation}
        \mathbf{r}[n] = \mathbf{r}(nT_s), \quad n \in \mathbb{Z}
    \end{equation}
    相应的，连续OC冲激响应矩阵$\mathbf{H}(t)$也被离散化为：
    \begin{equation}
        \mathbf{H}[n] = \mathbf{H}(nT_s), \quad n \in \mathbb{Z}
    \end{equation}
\end{definition}

\begin{remark}
    这里采用的是理想化的冲激响应采样模型。一般而言，仅令 $\mathbf{H}[n]=\mathbf{H}(nT_s)$ 并不必然得到与原连续时间系统完全等价的离散时间系统；真实工程实现还需要考虑保持器、抗混叠滤波器、数模转换器以及具体的离散化方法。后文采样定理讨论的是在理想带限与理想重构条件下的信息无损情形。
\end{remark}

根据存在性定理分析，无论连续离散还是离散时间 OC 变换的冲激响应矩阵$\mathbf{H}(t)$或$\mathbf{H}[n]$都是稳定的，存在
傅里叶变换。因此，可以使用频域方法分析采样过程。

\begin{theorem}[采样OC变换的频域关系]
    设连续时间OC变换在频域满足：
    \begin{equation}
        \tilde{\mathbf{Q}}(\mj\omega) = \tilde{\mathbf{H}}(\mj\omega) \tilde{\mathbf{R}}(\mj\omega)
    \end{equation}
    其中$\tilde{\mathbf{R}}(\mj\omega)$、$\tilde{\mathbf{H}}(\mj\omega)$、$\tilde{\mathbf{Q}}(\mj\omega)$分别是$\mathbf{r}(t)$、$\mathbf{H}(t)$、$\mathbf{q}(t)$的连续时间傅里叶变换。

    采样后的离散时间系统在频域满足：
    \begin{equation}
        \tilde{\mathbf{Q}}(\me^{\mj\varOmega}) = \dfrac{1}{T_s} \sum_{k=-\infty}^{\infty} \tilde{\mathbf{H}}\left(\dfrac{\varOmega + 2\pi k}{T_s}\right) \tilde{\mathbf{R}}\left(\dfrac{\varOmega + 2\pi k}{T_s}\right)
    \end{equation}
    其中$\varOmega = \omega T_s$是数字角频率。
\end{theorem}

\subsubsection{无失真采样条件}

\begin{theorem}[OC变换的无失真采样条件]
    要从离散OC序列$\mathbf{q}[n]$中无失真地恢复连续OC特征$\mathbf{q}(t)$，必须同时满足以下三个条件：

    \begin{enumerate}
        \item \textbf{现实描述带宽限制}：现实特征向量函数$\mathbf{r}(t)$是带限的，即存在$\omega_r^{\text{max}}$使得：
              \begin{equation}
                  \tilde{\mathbf{R}}(\mj\omega) = \mathbf{0}, \quad \forall |\omega| > \omega_r^{\text{max}}
              \end{equation}

        \item \textbf{系统带宽限制}：OC冲激响应矩阵$\mathbf{H}(t)$对应的传递函数矩阵是带限的，即存在$\omega_h^{\text{max}}$使得：
              \begin{equation}
                  \tilde{\mathbf{H}}(\mj\omega) = \mathbf{0}, \quad \forall |\omega| > \omega_h^{\text{max}}
              \end{equation}

        \item \textbf{采样率要求}：采样频率$\displaystyle\omega_s = \dfrac{2\pi}{T_s}$满足：
              \begin{equation}
                  \omega_s > 2 \cdot \min(\omega_r^{\text{max}}, \omega_h^{\text{max}})
                  \label{eq:nyquist-condition}
              \end{equation}
    \end{enumerate}
\end{theorem}

\begin{proof}
    要证明无失真恢复 $\mathbf{q}(t)$ 所需的最小采样率，需分析 $\mathbf{q}(t)$ 的带宽。

    在频域中，系统输出为：
    \begin{equation}
        \mathbf{Q}(\mj\omega) = \mathbf{H}(\mj\omega) \mathbf{R}(\mj\omega).
    \end{equation}
    已知：
    \begin{align*}
        \mathbf{R}(\mj\omega) & = \mathbf{0}, \quad \forall |\omega| > \omega_r^{\max}, \\
        \mathbf{H}(\mj\omega) & = \mathbf{0}, \quad \forall |\omega| > \omega_h^{\max}.
    \end{align*}

    考虑 $\mathbf{Q}(\mj\omega)$ 在 $|\omega| > \min(\omega_r^{\max}, \omega_h^{\max})$ 处的值。
    不失一般性，设 $\omega_r^{\max} \le \omega_h^{\max}$，则当 $|\omega| > \omega_r^{\max}$ 时，有 $\mathbf{R}(\mj\omega) = \mathbf{0}$。
    因此，对于 $|\omega| > \omega_r^{\max}$：
    \begin{equation}
        \mathbf{Q}(\mj\omega) = \mathbf{H}(\mj\omega) \cdot \mathbf{0} = \mathbf{0}.
    \end{equation}
    同理，若 $\omega_h^{\max} \le \omega_r^{\max}$，则当 $|\omega| > \omega_h^{\max}$ 时，$\mathbf{H}(\mj\omega) = \mathbf{0}$，同样导致 $\mathbf{Q}(\mj\omega) = \mathbf{0}$。

    故 $\mathbf{Q}(\mj\omega)$ 的带宽上限为：
    \begin{equation}
        B_q = \min(\omega_r^{\max}, \omega_h^{\max}).
    \end{equation}

    根据奈奎斯特-香农采样定理，要从采样序列 $\mathbf{q}[n] = \mathbf{q}(nT_s)$ 中无失真地恢复 $\mathbf{q}(t)$，采样频率须满足：
    \begin{equation}
        \omega_s > 2 B_q = 2 \min(\omega_r^{\max}, \omega_h^{\max}).
    \end{equation}
\end{proof}

\subsubsection{重构理论与理想插值}

\begin{theorem}[OC特征向量函数的重构]
    当满足上述无失真采样条件，并且已知双边采样序列时，OC特征向量函数可以从离散的OC特征向量序列完美重构：
    \begin{equation}
        \mathbf{q}(t) = \sum_{n=-\infty}^{\infty} \mathbf{q}[n] \cdot \text{sinc}\left(\dfrac{t - nT_s}{T_s}\right)
        \label{eq:ideal-reconstruction}
    \end{equation}
    其中$\text{sinc}(x) = \dfrac{\sin(\pi x)}{\pi x}$是理想插值函数。
\end{theorem}

\begin{proof}
    根据无失真采样条件，连续时间信号 $\mathbf{q}(t)$ 是带限的，其带宽 $B_q = \min(\omega_r^{\max}, \omega_h^{\max})$ 满足 $\omega_s = 2\pi/T_s > 2B_q$。

    考虑 $\mathbf{q}(t)$ 的连续时间傅里叶变换：
    \begin{equation}
        \tilde{\mathbf{Q}}(\mj\omega) = \int_{-\infty}^{\infty} \mathbf{q}(t) \me^{-\mj\omega t} \md t
    \end{equation}
    由带宽限制可知：
    \begin{equation}
        \tilde{\mathbf{Q}}(\mj\omega) = \mathbf{0}, \quad \forall |\omega| > \dfrac{\omega_s}{2}
    \end{equation}

    将 $\tilde{\mathbf{Q}}(\mj\omega)$ 在频域中以 $\omega_s$ 为周期进行延拓，得到周期谱：
    \begin{equation}
        \hat{\mathbf{Q}}(\mj\omega) = \sum_{k=-\infty}^{\infty} \tilde{\mathbf{Q}}(\mj(\omega - k \omega_s))
    \end{equation}
    由于 $\omega_s > 2B_q$，各周期副本互不重叠，无混叠现象。

    另一方面，采样序列 $\mathbf{q}[n] = \mathbf{q}(nT_s)$ 的离散时间傅里叶变换为：
    \begin{equation}
        \tilde{\mathbf{Q}}_d(\me^{\mj\omega T_s}) = \sum_{n=-\infty}^{\infty} \mathbf{q}[n] \me^{-\mj\omega nT_s}
    \end{equation}
    根据采样定理，离散谱与连续谱的关系为：
    \begin{equation}
        \tilde{\mathbf{Q}}_d(\me^{\mj\omega T_s}) = \dfrac{1}{T_s} \hat{\mathbf{Q}}(\mj\omega) = \dfrac{1}{T_s} \tilde{\mathbf{Q}}(\mj\omega), \quad |\omega| < \dfrac{\omega_s}{2}
    \end{equation}

    因此，在基带内可通过离散谱恢复连续谱：
    \begin{equation}
        \tilde{\mathbf{Q}}(\mj\omega) = T_s \cdot \tilde{\mathbf{Q}}_d(\me^{\mj\omega T_s}) \cdot \mathrm{rect}\left(\dfrac{\omega}{\omega_s}\right)
    \end{equation}
    其中 $\mathrm{rect}(u)$ 是矩形函数，在 $|u| < \dfrac{1}{2}$ 时为 1，否则为 0。

    对上述进行逆傅里叶变换。时域中，频域乘积对应于时域卷积：
    \begin{align*}
        \mathbf{q}(t) & = \mathcal{F}^{-1} \left\{ T_s \cdot \tilde{\mathbf{Q}}_d(\me^{\mj\omega T_s}) \cdot \mathrm{rect}\left(\dfrac{\omega}{\omega_s}\right) \right\}                             \\
                      & = \left( \sum_{n=-\infty}^{\infty} \mathbf{q}[n] \delta(t - nT_s) \right) * \left( T_s \cdot \dfrac{\omega_s}{2\pi} \cdot \text{sinc}\left(\dfrac{\omega_s t}{2\pi}\right) \right)
    \end{align*}
    其中 $T_s \cdot \dfrac{\omega_s}{2\pi} = T_s \cdot \dfrac{2\pi/T_s}{2\pi} = 1$，且 $\text{sinc}\left(\dfrac{\omega_s t}{2\pi}\right) = \text{sinc}\left(\dfrac{t}{T_s}\right)$。

    故：
    \begin{align*}
        \mathbf{q}(t) & = \sum_{n=-\infty}^{\infty} \mathbf{q}[n] \cdot \delta(t - nT_s) * \text{sinc}\left(\dfrac{t}{T_s}\right) \\
                      & = \sum_{n=-\infty}^{\infty} \mathbf{q}[n] \cdot \text{sinc}\left(\dfrac{t - nT_s}{T_s}\right)
    \end{align*}
    此即式\eqref{eq:ideal-reconstruction}。证明完毕。
\end{proof}

\begin{remark}
    若实际系统只保存 $n\geq 0$ 的因果截断序列，则式\eqref{eq:ideal-reconstruction} 的双边完美重构会变成近似重构；误差由截断长度、边界延拓方式和抗混叠滤波器共同决定。
\end{remark}

\subsubsection{采样定理在OC变换中的意义}

\begin{theorem}[离散到连续的完备性]
    如果采样过程满足香农条件，重构采用理想sinc插值，并且连续系统到离散系统的实现符合前述理想化采样模型，那么离散OC变换系统与连续OC变换系统在信息意义上是等价的：
    \begin{equation}
        \mathbf{q}(t) \xrightarrow{\text{采样}} \mathbf{q}[n] \xrightarrow{\text{重构}} \mathbf{q}(t)
    \end{equation}
    即采样和重构过程是信息无损的。
\end{theorem}

基于采样理论，连续时间OC变换可以在理想条件下通过离散时间系统无失真实现，只要：
\begin{enumerate}
    \item 选择合适的采样率$\omega_s$满足香农条件
    \item 设计适当的抗混叠滤波器
    \item 采用合适的插值方法进行重构
\end{enumerate}
这为自动OC变换机提供了理论基础。

%================= 快速 OC 变换 =================%
\subsection{快速 OC 变换}
\label{subsec:fast-oc}

从现实的角度考虑，任何角色不可能在无穷久远的时间就创立了，也不可能经过无穷大的时间仍然存在。在实际的自动OC变换机的使用中，可以设三者统一截断为长度 $L$ 的因果有限长序列：
\begin{equation}
    \mathbf{r}[n]\in\mathbb{R}^{N\times 1},\;
    \mathbf{H}[n]\in\mathbb{R}^{M\times N},\;
    \mathbf{q}[n]\in\mathbb{R}^{M\times 1},\qquad
    0\le n\le L-1.
\end{equation}
且 $L$ 为 2 的幂（不足则补零）。

此时，现实世界的设主要穿越到二次元世界成为其OC的过程主要包含以下步骤：
\begin{enumerate}
    \item 自动OC变换机的\textbf{采样模块}对设主的现实特征向量函数的采样与数字化
    \item 自动OC变换机的\textbf{正变换模块}对现实特征向量序列进行离散时间OC变换，得到OC特征向量序列
    \item 自动OC变换机的\textbf{重构模块}对OC特征向量序列重构，生成OC特征向量函数，从而转换成二次元OC角色
\end{enumerate}
当需要返回现实世界时，可知反向过程为：对OC特征向量函数采样，\textbf{反变换模块}进行离散时间OC反变换得到现实特征向量序列，进而重构生成现实特征向量函数

如果直接按定义计算卷积和，时间复杂度为 $O(NML^2)$，当 $N,M,L$ 较大时，计算量过于庞大，难以实时处理。为此，需引入快速算法以提升计算效率。
这就是OC变换在具体实现中最为重要的\textbf{快速OC变换}算法。

观察可知，离散时间OC变换的核心计算是 $N \times M$ 个逐元素卷积，每个卷积的长度为 $L$。
因此，提升自动OC变换机的计算效率，关键在于两点：
\begin{enumerate}
    \item \textbf{基于现代计算机缓存原理的矩阵乘法加速}：通过优化内存访问模式，提升矩阵乘法的效率
    \item \textbf{基于FFT的逐元素卷积加速}：利用快速傅里叶变换(FFT)将卷积计算从时域转到频域，大幅降低计算复杂度
\end{enumerate}

\subsubsection{基于现代计算机缓存原理的矩阵乘法加速}

矩阵乘法是离散时间OC变换的核心操作。由于直接计算矩阵乘法的时间复杂度为\(O(NM)\)（不考虑卷积的时间复杂度）
，已经相对较低，但在实际计算中，矩阵乘法的效率往往受限于内存访问模式和缓存命中率。

现代计算机通常采用分层缓存结构（L1、L2、L3缓存）来提高内存访问效率。通过优化矩阵乘法的内存访问模式，可以显著提升计算性能。

例如，下面这两段示例代码：
\begin{lstlisting}[language=python, caption=缓存不友好的离散时间OC变换的直接时域卷积]
def direct_convolution_bad_cache(M, N, L, H, r):
    q = np.zeros((M, L))
    # Cache-unfriendly access pattern
    for j in range(N):
        for i in range(M):
            conv = np.convolve(H[i, j], r[j], mode="full")[: L]
            q[i] += conv
    return q
\end{lstlisting}
\begin{lstlisting}[language=python, caption=缓存友好的离散时间OC变换的直接时域卷积]
def direct_convolution_good_cache(M, N, L, H, r):
    q = np.zeros((M, L))
    # Cache-friendly access pattern
    for i in range(M):
        temp = np.zeros(L)
        for j in range(N):
            conv = np.convolve(H[i, j], r[j], mode="full")[: L]
            temp += conv
        q[i] = temp
    return q
\end{lstlisting}

可以看出，第二段代码将内层循环分块处理，提升了缓存命中率，从而加速了OC变换中的矩阵乘法的计算。
实验表明，这种缓存优化策略在实际计算中可以带来一定的性能提升，尤其是在矩阵规模较大时。

\subsubsection{基于频域方法的卷积加速}

频域方法的核心在于利用离散傅里叶变换（DFT）的卷积定理，将计算复杂的时域卷积转化为计算简单的频域乘法，从而显著降低计算复杂度。这是降低OC变换计算负担的关键技术。

\begin{definition}[离散傅里叶变换（DFT）]
    对于一个长度为 $N_f$ 的离散序列 $x[n],\; n=0,1,\dots,N_f-1$，其\textbf{离散傅里叶变换（DFT） }$X[k]$ 定义为：
    \begin{equation}
        X[k] = \sum_{n=0}^{N_f-1} x[n] \cdot \me^{-\mj\frac{2\pi}{N_f}kn}, \quad k=0,1,\dots,N_f-1
    \end{equation}
    相应的\textbf{离散傅里叶逆变换（IDFT）}为：
    \begin{equation}
        x[n] = \dfrac{1}{N_f} \sum_{k=0}^{N_f-1} X[k] \cdot \me^{\mj\frac{2\pi}{N_f}kn}, \quad n=0,1,\dots,N_f-1
    \end{equation}
\end{definition}

对于进行了DFT后的序列，可以利用\textbf{循环卷积定理}来简化卷积计算。

\begin{theorem}[循环卷积定理]
    \label{thm:circular-conv}
    设序列 $\tilde{x}[n]$ 和 $\tilde{h}[n]$ 的长度均为 $N_f$，其循环卷积 $\tilde{y}[n]$ 定义为：
    \begin{equation}
        \tilde{y}[n] = \sum_{m=0}^{N_f-1} \tilde{h}[m] \tilde{x}[(n-m) \mod N_f]
    \end{equation}
    令 $\tilde{X}[k]$, $\tilde{H}[k]$, $\tilde{Y}[k]$ 分别为三者的DFT，则有：
    \begin{equation}
        \tilde{Y}[k] = \tilde{H}[k] \cdot \tilde{X}[k]
    \end{equation}
    即时域的循环卷积对应于频域的乘积。
\end{theorem}

卷积定理是连接时域与频域运算的桥梁，它指出：\textbf{时域中的循环卷积等价于频域中的点乘}。
然而，OC变换需要计算的是\textbf{线性卷积}，而非循环卷积。为了利用DFT计算两个长度为 $L$ 的因果序列 $h_{ij}[n]$ 和 $r_j[n]$ 的线性卷积，必须通过补零来避免混叠效应。

\begin{theorem}[基于DFT的线性卷积计算]
    \label{thm:dft-linear-conv}
    设 $h_{ij}[n]$ 和 $r_j[n]$ 均为长度 $L$ 的因果序列。将其各补零至长度 $N_f \ge 2L-1$（通常取 $N_f = 2L$），得到序列 $\tilde{h}_{ij}[n]$ 和 $\tilde{r}_j[n]$。此时，$\tilde{h}_{ij}[n]$ 与 $\tilde{r}_j[n]$ 的循环卷积的前 $L$ 个点，等于 $h_{ij}[n]$ 与 $r_j[n]$ 的线性卷积。
\end{theorem}

基于此定理，利用DFT计算线性卷积的步骤为：
\begin{enumerate}
    \item \textbf{补零}：对 $h_{ij}$ 和 $r_j$ 各补零至长度 $N_f = 2L$。
    \item \textbf{正变换}：计算两者的DFT：$H_{ij}[k] = \mathrm{DFT}(\tilde{h}_{ij})$, $R_j[k] = \mathrm{DFT}(\tilde{r}_j)$。
    \item \textbf{频域相乘}：$Y_{ij}[k] = H_{ij}[k] \cdot R_j[k]$。
    \item \textbf{反变换}：计算IDFT并取前 $L$ 个点：$(h_{ij} * r_j)[n] = \mathrm{IDFT}(Y_{ij})[n],\; n=0,\dots,L-1$。
\end{enumerate}

直接计算一个 $N_f$ 点DFT的时间复杂度为 $O(N_f^2)$。若采用此方法计算卷积，总复杂度为 $O(N_f^2) = O(L^2)$，与直接时域卷积相当，并未带来优势。
为此，引入快速傅里叶变换（FFT）算法以提升DFT的计算效率。

\textbf{快速傅里叶变换（FFT）}是一类高效计算DFT的算法总称，其核心思想是分治策略，通过将大的DFT分解为多个小DFT来减少计算量。%最著名的Cooley-Tukey算法将 $N_f$（$N_f$ 为2的幂）点DFT分解为两个 $N_f/2$ 点DFT，并递归进行。

FFT有多种实现方式，其中最为经典的是Cooley-Tukey算法。当变换点数 $N_f$ 是2的幂时，可采用基2算法，其主要分为按时间抽选(DIT)和按频率抽选(DIF)两种。

\begin{algorithm}
    \caption{基2 FFT算法（按时间抽选, DIT）}
    \label{alg:fft-dit}
    \begin{algorithmic}[1]
        \Require 输入序列 $x[n]$, $n=0,1,\dots,N_f-1$，$N_f$ 为2的幂
        \Ensure 频域序列 $X[k] = \mathrm{FFT}\{x[n]\}$，$k=0,1,\dots,N_f-1$
        \State 初始化：将输入序列按比特反序排列
        \For {$s = 1$ to $\log_2 N_f$} \Comment{级数循环}
        \State $m \gets 2^s$ \Comment{当前级子DFT长度}
        \State $W_m \gets \me^{-\mj\frac{2\pi}{m}}$ \Comment{旋转因子基}
        \For {$i = 0$ to $m/2-1$}
        \State $W_m^i \gets \me^{-\mj\frac{2\pi i}{m}}$ \Comment{计算旋转因子}
        \EndFor
        \For {$k = 0$ to $N_f-1$ step $m$} \Comment{蝶形组循环}
        \For {$l = 0$ to $m/2-1$} \Comment{蝶形运算循环}
        \State $u \gets x[k + l]$
        \State $v \gets W_m^l \cdot x[k + l + m/2]$
        \State $x[k + l] \gets u + v$
        \State $x[k + l + m/2] \gets u - v$
        \EndFor
        \EndFor
        \EndFor
        \State \Return $x$ \Comment{此时 $x$ 中存储的即为按自然顺序排列的 $X[k]$}
    \end{algorithmic}
\end{algorithm}

算法\ref{alg:fft-dit}的核心思想是：将输入序列按奇偶索引分解为两个子序列，递归计算子序列的DFT，再通过蝶形运算组合得到最终结果。其名称“按时间抽选”源于此分解是在时域（即输入序列）上进行的。

\begin{algorithm}[H]
    \caption{基2 FFT算法（按频率抽选, DIF）}
    \label{alg:fft-dif}
    \begin{algorithmic}[1]
        \Require 输入序列 $x[n]$, $n=0,1,\dots,N_f-1$，$N_f$ 为2的幂
        \Ensure 频域序列 $X[k] = \mathrm{FFT}\{x[n]\}$，$k=0,1,\dots,N_f-1$
        \For {$s = \log_2 N_f$ down to $1$} \Comment{级数循环}
        \State $m \gets 2^s$ \Comment{当前级子DFT长度}
        \State $W_m \gets \me^{-\mj\frac{2\pi}{m}}$ \Comment{旋转因子基}
        \For {$i = 0$ to $m/2-1$}
        \State $W_m^i \gets \me^{-\mj\frac{2\pi i}{m}}$ \Comment{计算旋转因子}
        \EndFor
        \For {$k = 0$ to $N_f-1$ step $m$} \Comment{蝶形组循环}
        \For {$l = 0$ to $m/2-1$} \Comment{蝶形运算循环}
        \State $u \gets x[k + l]$
        \State $v \gets x[k + l + m/2]$
        \State $x[k + l] \gets u + v$
        \State $x[k + l + m/2] \gets (u - v) \cdot W_m^l$
        \EndFor
        \EndFor
        \EndFor
        \State 将输出序列按比特反序排列
        \State \Return $x$ \Comment{此时 $x$ 中存储的即为按比特反序排列的 $X[k]$}
    \end{algorithmic}
\end{algorithm}

算法\ref{alg:fft-dif}与DIT算法的主要区别在于：
\begin{itemize}
    \item \textbf{分解方向}：DIF首先在频域上将DFT输出序列按奇偶索引分解，故称“按频率抽选”。
    \item \textbf{蝶形运算}：在DIF中，旋转因子乘法位于蝶形运算的后半支路。
    \item \textbf{顺序要求}：DIT要求输入为比特反序，输出为自然顺序；DIF则要求输入为自然顺序，输出为比特反序。
\end{itemize}

两种算法均将计算复杂度从 $O(N_f^2)$ 降至 $O(N_f \log N_f)$，且具有相同的计算量和数值精度，在实际应用中可根据具体需求选择。

\begin{theorem}[FFT加速的线性卷积计算复杂度]
    \label{thm:fft-element}
    设 $h_{ij}[n]$ 和 $r_j[n]$ 均为长度 $L$ 的因果序列（不足则补零至长度 $N_f=2L$，且 $N_f$ 为2的幂）。利用FFT计算线性卷积的复杂度为：
    \begin{equation}
        O(N_f \log N_f) = O(L \log L)
    \end{equation}
    相比之下，直接时域卷积的复杂度为 $O(L^2)$。
\end{theorem}

\begin{proof}
    FFT算法将DFT的计算复杂度从 $O(N_f^2)$ 降至 $O(N_f \log N_f)$。在上述卷积计算步骤中：
    \begin{enumerate}
        \item 两次FFT（计算 $H_{ij}[k]$ 和 $R_j[k]$）的复杂度为 $2 \times O(N_f \log N_f)$。
        \item 一次频域点乘的复杂度为 $O(N_f)$。
        \item 一次IFFT的复杂度为 $O(N_f \log N_f)$。
    \end{enumerate}
    故总复杂度为 $O(N_f \log N_f) = O(L \log L)$。
\end{proof}

在OC变换的完整计算中，需要对所有 $N \times M$ 个通道分别计算卷积。若采用直接时域卷积，总复杂度为 $O(NML^2)$；而采用FFT加速后，总复杂度降至：
\begin{equation}
    O(NM \cdot L \log L)
\end{equation}

值得注意的是，在实际实现中可以利用变换域的系数复用进一步优化。具体而言：
\begin{itemize}
    \item \textbf{系数预计算}：离散OC冲激响应矩阵 $\mathbf{H}$ 通常是固定的系统参数，可以预先计算并存储所有 $H_{ij} = \mathrm{FFT}(h_{ij})$。
    \item \textbf{输入变换复用}：对于每个输入通道 $r_j$，只需计算一次 $\mathrm{FFT}(r_j)$，即可复用于所有 $M$ 个输出通道的计算。
    \item \textbf{并行计算}：各输出通道 $q_i$ 的计算相互独立，适合并行处理。
\end{itemize}

优化后的计算流程如算法\ref{alg:fast-oc-transform}所示。
\begin{algorithm}
    \caption{基于FFT的快速OC变换算法}
    \label{alg:fast-oc-transform}
    \begin{algorithmic}[1]
        \State \textbf{预计算：} 对所有 $i,j$，计算 $H_{ij} \gets \mathrm{FFT}(h_{ij})$ 并存储
        \For{$j = 1$ to $N$}
        \State $R_j \gets \mathrm{FFT}(r_j)$ \Comment{输入变换}
        \EndFor
        \For{$i = 1$ to $M$}
        \State 初始化 $Q_i \gets 0$ \Comment{频域累加器}
        \For{$j = 1$ to $N$}
        \State $Q_i \gets Q_i + H_{ij} \odot R_j$ \Comment{频域乘积累加}
        \EndFor
        \State $q_i \gets \mathrm{IFFT}(Q_i)$ \Comment{输出变换}
        \EndFor
    \end{algorithmic}
\end{algorithm}

该算法将总复杂度从 $O(NML^2)$ 显著降低至 $O(NML \log L)$，在 $L$ 较大时带来数个数量级的加速比，使得实时OC变换成为可能。

%================= 综合算法 =================%
\subsubsection{综合优化算法：缓存优化与FFT结合}

在实际应用中，为了最大化OC变换的计算效率，可以将缓存优化策略与FFT加速方法相结合，形成综合优化算法。该算法的核心思想是：

\begin{enumerate}
    \item \textbf{频域计算主体框架}：采用基于FFT的频域卷积作为主要计算范式，将时间复杂度从 $O(NML^2)$ 降低至 $O(NML \log L)$
    \item \textbf{缓存优化的内存访问}：在矩阵乘法累加阶段，优化内存访问模式，提升缓存命中率，减少内存带宽瓶颈
    \item \textbf{并行计算架构}：利用现代多核处理器的并行计算能力，同时对多个输出通道或输入通道进行计算
\end{enumerate}

该综合算法在保持FFT算法低时间复杂度的同时，通过优化内存访问模式进一步减少了常数因子，在实际硬件上能够达到接近理论峰值性能的计算效率。

%================= 算法对比表格 =================%
\subsubsection{算法性能对比分析}

表\ref{tab:algorithm-comparison} 详细对比了五种OC变换算法的性能特征。

\begin{table}[htbp]
    \centering
    \caption{OC变换算法性能对比}
    \label{tab:algorithm-comparison}
    \begin{tabular}{lccc}
        \toprule
        \textbf{算法类型} & \textbf{时间复杂度}  & \textbf{空间复杂度} & \textbf{速度} \\
        \midrule
        缓存不友好的时域卷积    & $O(NML^2)$      & $O(NML)$       & 慢           \\
        缓存友好的时域卷积     & $O(NML^2)$      & $O(NML)$       & 较慢          \\
        基于原始DFT算法     & $O(NML^2)$      & $O(NML)$       & 最慢          \\
        基于FFT的频域算法    & $O(NML \log L)$ & $O(NML)$       & 快           \\
        综合优化算法        & $O(NML \log L)$ & $O(NML)$       & 很快          \\
        \bottomrule
    \end{tabular}
\end{table}

\begin{remark}
    有趣的是，基于原始DFT算法的“快速”OC变换算法其实一点也不快。这是由于DFT的计算复杂度为 $O(N_f^2)=O(L^2)$，与直接时域卷积的 $O(L^2)$ 相当，且DFT的常数因子较大，导致其在实际应用中反而更慢。
\end{remark}

% 值得注意的是，当 $L$ 较小时，FFT的常数因子可能使得频域方法反而不如优化后的时域方法高效。因此在实际应用中，通常设置一个阈值，根据问题规模动态选择最优算法。

\subsubsection{数值实验与性能分析}

% \subsection{快速OC变换算法性能评估}

附录C给出了\hyperref[chap:fast_oc_transform]{快速OC变换算法示例：Python代码}。经过系统测试（均关闭了多线程并行加速），在不同规模问题下各种算法的性能对比如表\ref{tab:performance-comparison-new}所示，对应的运行时间与加速比曲线如图\ref{fig:performance-comparison}所示。

\begin{table}[htbp]
    \centering
    \caption{快速OC变换算法性能比较}
    \label{tab:performance-comparison-new}
    \begin{tabular}{lcccc}

        \textbf{测试规模} & \textbf{算法名称}                    & \textbf{运行时间 (s)} & \textbf{加速比} & \textbf{平均绝对误差} \\
        \midrule
        \multirow{5}{*}{\(N=4, L=32, M=2\)}
                      & direct\_convolution\_bad\_cache  & 0.000062          & 1.000000     & 0.000000e+00    \\
                      & direct\_convolution\_good\_cache & 0.000026          & 2.425926     & 0.000000e+00    \\
                      & dft\_direct                      & 0.000469          & 0.133198     & 2.942432e-05    \\
                      & fft\_bad\_cache                  & 0.011228          & 0.005563     & 1.081682e-17    \\
                      & fft\_good\_cache                 & 0.000048          & 1.290640     & 1.081682e-17    \\
        \midrule
        \multirow{5}{*}{\(N=8, L=128, M=4\)}
                      & direct\_convolution\_bad\_cache  & 0.000321          & 1.000000     & 0.000000e+00    \\
                      & direct\_convolution\_good\_cache & 0.000299          & 1.072510     & 0.000000e+00    \\
                      & dft\_direct                      & 0.005594          & 0.057364     & 1.162342e-14    \\
                      & fft\_bad\_cache                  & 0.000271          & 1.183817     & 5.231555e-18    \\
                      & fft\_good\_cache                 & 0.000107          & 2.997773     & 5.231555e-18    \\
        \midrule
        \multirow{5}{*}{\(N=16, L=512, M=8\)}
                      & direct\_convolution\_bad\_cache  & 0.005002          & 1.000000     & 0.000000e+00    \\
                      & direct\_convolution\_good\_cache & 0.003951          & 1.266188     & 0.000000e+00    \\
                      & dft\_direct                      & 0.251821          & 0.019865     & 8.289528e-16    \\
                      & fft\_bad\_cache                  & 0.001900          & 2.632292     & 3.268654e-18    \\
                      & fft\_good\_cache                 & 0.001729          & 2.892872     & 3.268654e-18    \\
        \midrule
        \multirow{5}{*}{\(N=32, L=1024, M=16\)}
                      & direct\_convolution\_bad\_cache  & 0.051716          & 1.000000     & 0.000000e+00    \\
                      & direct\_convolution\_good\_cache & 0.060097          & 0.860547     & 0.000000e+00    \\
                      & dft\_direct                      & 3.019521          & 0.017127     & 6.469994e-16    \\
                      & fft\_bad\_cache                  & 0.014311          & 3.613711     & 2.282972e-18    \\
                      & fft\_good\_cache                 & 0.014950          & 3.459262     & 2.282972e-18    \\
        \midrule
        \multirow{5}{*}{\(N=64, L=2048, M=32\)}
                      & direct\_convolution\_bad\_cache  & 3.470409          & 1.000000     & 0.000000e+00    \\
                      & direct\_convolution\_good\_cache & 3.476582          & 0.998225     & 0.000000e+00    \\
                      & dft\_direct                      & 44.916513         & 0.077264     & 4.598533e-17    \\
                      & fft\_bad\_cache                  & 0.115147          & 30.139038    & 8.720734e-19    \\
                      & fft\_good\_cache                 & 0.111049          & 31.251235    & 8.720734e-19    \\
        \bottomrule
    \end{tabular}
\end{table}

\begin{figure}[htbp]
    \centering
    \includegraphics[width=0.8\textwidth]{../figure/oc_transform_benchmark.png}
    \caption{不同规模下各OC变换算法运行时间对比}
    \label{fig:performance-comparison}
\end{figure}

从实验结果可以得出以下重要结论：

\begin{enumerate}
    \item \textbf{计算效率层次分明}：算法性能呈现明显的规模依赖性。在小规模问题（$N=4,L=32$）中，直接卷积优化版本表现最佳；而在大规模问题（$N=64,L=2048$）中，FFT方法实现了30倍以上的加速比，充分体现了其$O(MNL\log L)$复杂度的优势。

    \item \textbf{缓存优化效果显著}：对比"good\_cache"与"bad\_cache"版本，优化的内存访问模式在不同规模下带来稳定性能提升。在$N=4,L=32$情况下，direct\_convolution\_good\_cache比direct\_convolution\_bad\_cache快2.43倍；在$N=8,L=128$时，fft\_good\_cache比fft\_bad\_cache快2.53倍，凸显了缓存友好算法设计的重要性。

    \item \textbf{数值精度卓越}：所有FFT方法的平均绝对误差均在$10^{-17}$-$10^{-19}$量级，展现了极高的数值稳定性。DFT直接法由于数值累积误差，精度相对较低（$10^{-5}$-$10^{-17}$）但仍满足工程应用需求。直接卷积方法由于是精确计算，误差为零。

    \item \textbf{算法适用场景分析}：
          \begin{itemize}
              \item \textbf{缓存友好的直接卷积}：适合小规模问题（$L \leq 512$），实现简单且性能优秀
              \item \textbf{FFT方法}：在大规模问题（$L \geq 1024$）中优势明显，特别适合实时OC变换需求，在$N=64,L=2048$时达到31倍加速比，这是目前自动OC变换机的首选算法
              \item \textbf{DFT直接法}：计算复杂度$O(MNL^2)$导致性能最差，在$N=64$时耗时44.9秒，仅适用于理论验证
          \end{itemize}

    \item \textbf{系统稳定性验证}：所有测试案例中最大绝对和为0.909091，远小于1，满足绝对可和条件，确保了OC变换系统的BIBO稳定性。输出范围在不同算法间保持一致，进一步验证了系统的数值稳定性。

    \item \textbf{复杂度理论验证}：实测结果与理论分析高度一致。当$L$从32增长到2048（64倍），直接卷积计算时间从0.000062s增长到3.470409s（约55970倍），接近$O(L^2)$的4096倍理论增长；而FFT方法从0.000048s增长到0.111049s（约2313倍），接近$O(L\log L)$的理论增长趋势。
\end{enumerate}

在实际工程实现中，FFT方法的性能对缓存利用率和数据布局极为敏感。测试表明，通过优化数据访问模式（如fft\_good\_cache），在不同规模问题上可获得稳定性能提升。这对自动OC变换机的设计者提出了重要启示：需要根据问题规模动态选择最优算法，并在算法实现中充分考虑现代处理器的内存层次结构特性。

\section{OC的数字图像生成}\label{sec:ocimage}
本节讨论 OC 角色图像是如何从文字设定逐步变成可反复使用的视觉形象。这里的“数字图像生成”不只指自动生成图像，也包括画师使用数位板、绘图软件、图层、笔刷和参考资料绘制 OC 的过程。

在前文的 OC 变换理论中，OC 的性质被抽象为特征向量函数 $\mathbf{q}(t)$。但在实际约稿或绘制中，设主不会把 $\mathbf{q}(t)$ 以数学表格交给画师。画师真正接收到的通常是文字描述、设定图、色板、参考图、剧情说明或性格说明。

因此，OC 的数字图像生成可以分成两个主要方向：
\begin{enumerate}
    \item \textbf{文字设定到设定图}：把语言描述转化为第一套稳定的视觉基准。
    \item \textbf{设定图到更多稿件}：在已有视觉基准上，绘制不同姿势、表情、服装状态、场景和构图中的 OC 图像。
\end{enumerate}

这两个方向对应了 OC 图像生产中的两个关键问题：首先要让 OC “长什么样”被确定下来，然后才能让这个 OC 在更多画面中保持一致。

\subsection{基于文字设定的图像生成}\label{subsec:text_based_gen}

本小节讨论第一条链路：\textbf{文字设定到设定图}。文字设定是最常见的 OC 初始信息载体。例如设主可能写下：“银色长发，红色眼睛，黑色连衣裙，性格冷淡，使用冰系能力。”这些描述能够告诉画师角色的大致方向，但它们还不是可直接复用的视觉基准。

文字表达的是概念，设定图必须给出具体形状。同样是“银色长发”，可以是直发、卷发、低马尾、双马尾、齐腰长发，也可以是带蓝灰色阴影的冷银色，或接近白色的高亮银色。同样是“黑色连衣裙”，可以是哥特裙、礼服裙、制服裙、战斗服式短裙，也可以带蕾丝、金属扣、披肩或不规则裙摆。文字越短，第一张设定图需要补全的空白就越多。

在特征向量函数的语言下，文字设定可以看作对 $\mathbf{q}(t)$ 的低密度描述。它只约束了部分特征，例如发色、瞳色、服装类别、性格气质和能力方向，却没有完全规定发型轮廓、服装层级、材质、色彩比例和细节位置。

这种空白主要来自三类信息缺失：
\begin{enumerate}
    \item \textbf{形状缺失}：文字说出了“有什么”，但没有说明轮廓、比例和结构。例如“披风”可以很短，也可以拖地；“角”可以是羊角、龙角或鹿角。
    \item \textbf{颜色与材质缺失}：文字说出了颜色名称，但没有说明明度、饱和度、纹理和反光方式。例如“白色外套”可以是棉布、皮革、羽绒或金属装甲。
    \item \textbf{关系缺失}：文字列出了多个元素，但没有说明它们如何组合。例如“黑裙、红宝石、银链”并不自动说明宝石在胸口、腰带还是发饰上。
\end{enumerate}

因此，文字转设定图本质上是一个“解释并定型”的过程。画师先读懂文字，再根据经验、审美和画风，把没有被明说的部分补出来，并把补全结果固定为可供后续参考的视觉基准。

这一过程可用简化模型表示为：
\begin{equation}
    I_{\text{设定图}}=G_{\text{text}\to\text{ref}}(T,A),
\end{equation}
其中 $T$ 表示文字设定，$A$ 表示画师的解释、经验与画风，$I_{\text{设定图}}$ 表示生成后的设定图。这个公式不是用于计算，而是说明：设定图不是由文字单独决定的，它还包含画师对文字空白处的视觉判断。

\begin{example}[文字设定的多解性]
    设定“白发、蓝眼、魔法师、安静”至少可以生成以下几种完全不同的图像方向：
    \begin{itemize}
        \item 穿长袍、手持法杖的古典魔法师；
        \item 穿学院制服、拿魔导书的学生型角色；
        \item 穿现代外套、使用冰晶能力的都市幻想角色；
        \item 头戴面纱、带宗教感饰品的神秘祭司。
    \end{itemize}
    这些结果都没有违背文字设定，但它们会形成不同的设定图。
\end{example}

文字转设定图完成后，OC 会获得第一套稳定的视觉基准。此时的设定图可以近似看作 $\mathbf{q}(t)$ 在某个初始状态 $t_0$ 下的视觉采样：
\begin{equation}
    I_{\text{设定图}}\approx \operatorname{Sample}(\mathbf{q}(t_0))
\end{equation}
这里的“采样”表示把抽象的 OC 特征转化为一张具体可见的图像。设定图无法包含 OC 的全部可能状态，但它固定了后续稿件最重要的参考基准。

\subsection{基于特征向量函数的图像生成}\label{subsec:feature_based_gen}

本小节讨论第二条链路：\textbf{设定图到更多稿件}。当设定图已经存在时，后续稿件不再从零开始解释文字，而是以设定图为视觉基准，绘制 OC 在不同动作、表情、场景和画风要求下的图像。

设定图不是普通参考图，而是 OC 特征的视觉说明书。它把原本需要大量文字解释的内容，压缩到可观察的形状、颜色和结构中。更多稿件的任务，是在不破坏这些稳定特征的前提下，让 OC 出现在新的画面中。

\subsubsection{设定图提供的稳定特征}

一张完整的设定图通常不只是“画得好看的一张图”。它承担的是说明功能。不同部分解决不同问题：
\begin{enumerate}
    \item \textbf{正面图}：说明角色的主要识别点，如发型、脸型、服装正面结构、主要配色。
    \item \textbf{侧面图与背面图}：说明从其他角度看不到的结构，如发尾长度、披风连接方式、背部装饰、武器挂载位置。
    \item \textbf{色板}：固定颜色，避免“红色”“蓝色”“银色”等词在不同画师那里产生偏差。
    \item \textbf{服装拆解}：说明衣服的层级关系。例如外套在披肩外面还是里面，腰带压住裙子还是挂在外层。
    \item \textbf{表情与姿态参考}：说明角色的性格如何体现在脸部和身体语言中。
    \item \textbf{细节放大}：说明眼睛纹样、徽章、耳饰、武器纹路等小尺寸但高识别度的部分。
\end{enumerate}

这些信息共同降低了后续稿件的偏差。纯文字需要画师想象“它长什么样”，而设定图直接告诉画师“它大致就长这样”。所以设定图越完整，后续稿件越容易保持角色一致性。

\subsubsection{从设定图到新稿件}

画师看到设定图后，通常不会简单地把原图照抄到新画面中。真正的绘制过程更接近“提取稳定特征，再在新画面中重组”。例如要画一个新动作，画师需要先判断哪些东西必须保留，哪些东西可以随着动作变化。

这个理解过程可以分成几步：
\begin{enumerate}
    \item \textbf{读取轮廓}：先确认角色的大形。例如头发外轮廓、服装剪影、武器长度、翅膀或尾巴的位置。
    \item \textbf{确认比例}：判断头身比、肩宽、腿长、手脚大小等。比例一变，角色气质也会改变。
    \item \textbf{分解层级}：理解衣服和饰品的前后关系。比如领结压在衬衫上，披风压在外套后，腰封压住裙腰。
    \item \textbf{识别重点}：找出这个 OC 最不能丢的识别点。可能是眼睛形状、发饰、角、配色、纹身、武器或某个特殊符号。
    \item \textbf{推断变化}：当角色转身、抬手、奔跑或受风吹时，头发、布料、饰品会如何移动。
    \item \textbf{匹配画风}：把设定图转化为当前稿件的画风。厚涂、赛璐璐、Q版、立绘和头像对细节的取舍不同。
\end{enumerate}

设定图给出的是约束，画师要在这些约束内完成新的图像。好的后续稿件不是机械复制设定图，而是在新构图中仍然保留“这是同一个 OC”的识别性。

这一过程可用简化模型表示为：
\begin{equation}
    I_{\text{稿件}}=G_{\text{ref}\to\text{art}}(I_{\text{设定图}},C),
\end{equation}
其中 $I_{\text{设定图}}$ 是视觉基准，$C$ 表示新稿件的构图、动作、表情、光照、场景和画风要求，$I_{\text{稿件}}$ 是最终稿件。

若使用特征向量函数描述，则后续稿件可以近似看作 $\mathbf{q}(t)$ 在另一状态 $t_1$ 下的视觉采样：
\begin{equation}
    I_{\text{稿件}}\approx \operatorname{Sample}(\mathbf{q}(t_1),C)
\end{equation}
其中 $t_1$ 不一定表示真实时间，也可以表示新的表情状态、动作状态或剧情状态。

\subsection{数字绘图中的生成流程}

从数字绘图角度看，一张 OC 稿件通常经过以下阶段：
\begin{enumerate}
    \item \textbf{草图阶段}：确定构图、动作和大比例。此时重点不是细节，而是角色是否站得住、姿势是否符合性格。
    \item \textbf{结构整理}：修正人体、透视、服装层级和道具位置。复杂 OC 往往在这一阶段最容易出错。
    \item \textbf{线稿阶段}：把结构固定下来，决定哪些线条是轮廓线，哪些线条是内部细节。
    \item \textbf{固有色阶段}：铺上角色的基础颜色，并检查是否符合色板或设主要求。
    \item \textbf{明暗阶段}：加入阴影和光源，让角色从平面设定变成有体积的图像。
    \item \textbf{材质与特效阶段}：表现布料、金属、皮革、宝石、火焰、冰晶、魔法阵等特殊效果。
    \item \textbf{校对阶段}：检查发型、配色、饰品、服装层级和角色气质是否偏离原设。
\end{enumerate}

在这个流程中，越靠前的阶段越影响整体，越靠后的阶段越影响完成度。若草图阶段已经误解了角色比例，后面即使上色精美，也可能不像原来的 OC。若结构正确但材质处理不足，角色仍然能被认出，只是完成度较低。

\subsection{两类生成误差}

OC 图像生成中的偏差主要有两类。

第一类是\textbf{文字到设定图的误差}。它来自语言描述的不完整。例如设主写了“短外套”，但没有说明长度、材质、袖口、扣子和内搭，画师就必须做出选择。这个阶段的误差会影响 OC 的第一套视觉基准。

第二类是\textbf{设定图到更多稿件的误差}。它来自视角、动作、表情、光影、构图和画风转换。例如设定图中正面可见的胸口徽章，在侧身动作中可能被手臂遮挡；长发在奔跑时会飘起；Q版稿件会省略部分复杂装饰。这个阶段的重点不是逐像素一致，而是保持核心识别点一致。

因此，文字、设定图和更多稿件之间不是简单复制关系，而是逐层约束关系。文字给出方向，设定图固定外观，更多稿件在此基础上扩展 OC 的可见状态。

\subsection{本节小结}

OC 的数字图像生成主要包括两个方向：文字设定转设定图，设定图转更多稿件。前者解决“这个 OC 的视觉基准是什么”，后者解决“这个 OC 如何在不同画面中继续保持一致”。

在 OC 变换理论的语言中，文字设定是对 $\mathbf{q}(t)$ 的低密度描述，设定图是对 $\mathbf{q}(t_0)$ 的初始视觉采样，更多稿件则是在不同状态和绘制条件下对 $\mathbf{q}(t)$ 的再次采样。这样，文字、设定图和稿件就被统一到同一套 OC 特征描述框架中。
